% !TEX program = pdflatex
\documentclass{article}

% ============================================================
%  NeurIPS 2026 Style
% ============================================================
\usepackage[final]{neurips_2026}   % camera-ready; use [preprint] for arXiv

% ============================================================
%  Packages
% ============================================================
\usepackage[T1]{fontenc}
\usepackage[utf8]{inputenc}
\usepackage{amsmath,amssymb,amsfonts,amsthm}
\usepackage{algorithmic}
\usepackage{algorithm}
\usepackage{booktabs}
\usepackage{multirow}
\usepackage{graphicx}
\usepackage{xcolor}
\usepackage{url}
\usepackage{hyperref}
\usepackage{bbm}
\usepackage{bm}
\usepackage{subcaption}
\usepackage{enumitem}
\usepackage{natbib}
\usepackage{microtype}
\usepackage{cleveref}

\hypersetup{
  colorlinks=true,
  linkcolor=blue!60!black,
  citecolor=blue!60!black,
  urlcolor=blue!70!black
}

\setlist{nosep}

% ============================================================
%  Theorem Environments
% ============================================================
\newtheorem{theorem}{Theorem}
\newtheorem{lemma}[theorem]{Lemma}
\newtheorem{corollary}[theorem]{Corollary}
\newtheorem{proposition}[theorem]{Proposition}
\newtheorem{definition}[theorem]{Definition}
\newtheorem{remark}[theorem]{Remark}

% ============================================================
%  Custom Commands
% ============================================================
\newcommand{\ie}{\textit{i.e.}}
\newcommand{\eg}{\textit{e.g.}}
\newcommand{\etal}{\textit{et al.}}
\newcommand{\wrt}{w.r.t.}
\newcommand{\UDL}{\mathrm{UDL}}
\newcommand{\R}{\mathbb{R}}
\newcommand{\E}{\mathbb{E}}
\newcommand{\Nset}{\mathcal{S}}  % Normal set (avoid \mathcal{N} ≈ Gaussian confusion)
\newcommand{\T}{\mathcal{T}}
\newcommand{\calR}{\mathcal{R}}
\newcommand{\norm}[1]{\left\lVert#1\right\rVert}
\newcommand{\abs}[1]{\left\lvert#1\right\rvert}
\DeclareMathOperator{\erf}{erf}
\DeclareMathOperator{\tr}{tr}
\DeclareMathOperator{\rank}{rank}
\DeclareMathOperator{\diag}{diag}
\DeclareMathOperator*{\argmax}{arg\,max}
\DeclareMathOperator*{\argmin}{arg\,min}

% ============================================================
%  Title
% ============================================================
\title{The Universal Deviation Principle: Multi-Spectrum Sensitivity Guarantees\\for Anomaly Detection}

\author{
  Odeyemi Olusegun Israel \\
  Independent Researcher \\
  \texttt{segettii@gmail.com} \\
  ORCID: \href{https://orcid.org/0009-0004-2309-6436}{0009-0004-2309-6436}
}

\begin{document}
\maketitle

% ============================================================
%  Abstract
% ============================================================
\begin{abstract}
We introduce the \textbf{Universal Deviation Principle (UDP)}, a theoretically grounded framework for multi-perspective anomaly detection. The core principle is that, under explicit assumptions on operator-family richness and anomaly-normal separation in feature space, a point which is \emph{meaningfully different from the normal set} cannot simultaneously match normality under all structural perspectives. We formalise this via \emph{deviation-separating operator families}, collections of differentiable maps $\{\Phi_k\}_{k=1}^K$ whose stacked Jacobian has full column rank, yielding local distinguishability around a reference normal state $\bm{\mu}$. To connect the sensitivity theory to detection, we define a normal set $\Nset$ and a deviation score based on distance of $\Phi(\mathbf{x})$ to $\Phi(\Nset)$; under an explicit separation condition between normal and anomaly sets in operator space, at least one operator deviates by a positive margin. All theoretical results are conditional on three stated assumptions (observability, operator-family richness, calibration discipline). We prove five results: (1)~local distinguishability via the Inverse Function Theorem; (2)~global sensitivity via compactness and the Extreme Value Theorem with a uniform bound $\underline{\sigma} > 0$; (3)~finite-sample guarantees via Hoeffding concentration; (4)~fusion weights maximising worst-case directional sensitivity via a semidefinite programme; and (5)~a greedy operator selection algorithm with $(1-1/e)$-approximation guarantee via monotone submodularity. We instantiate the framework with five spectrum projection operators (statistical, dynamical, spectral, geometric, exponential) and introduce a \textbf{Boundary-Centred Dimension Magnifier} paired with QDA that, given labels, achieves strong supervised performance. The unsupervised MDN score provides a label-free alternative (mean AUC $0.829$ across five benchmarks).

We present three \textbf{System Mode engines}---MolecularEngine (Lennard-Jones), GravityModeEngine (Gaussian attraction with Coulomb repulsion), and HybridGravityEngine (adaptive blend)---each stabilised by a \textbf{Lyapunov~v2 controller} that provides three convergence guarantees: (i)~Armijo descent with input-to-state stability (ISS) margin for treating pairwise forces as bounded perturbations, (ii)~La~Salle invariance convergence to the largest invariant set where $\nabla E = 0$, and (iii)~$C^1$-smooth barrier-augmented force clamping that eliminates chattering. Prediction uses a \textbf{fused topological scoring architecture}---Morse topology alarm, Betti barcode suite, and UDL post-simulation operators---that is structurally decoupled from the solver Lyapunov signal, ensuring 100\% topologically stable prediction. A \textbf{false-alarm-rate (FAR) calibrator} enables regulatory-compliant deployment.

Beyond the 5~base operators, we extend the framework to 14~composable spectrum operators across five mathematical families. A correlation audit across all 91~operator pairs reveals only one redundant pair (Wavelet$\leftrightarrow$Phase, $\rho{=}0.92$), validating the multi-spectrum thesis. We introduce \textbf{top-$k$ coverage} as a complementary metric to AUC. UDL-Hybrid with auto-selected projection strategy achieves mean AUC $0.972$ and 91\% mean coverage across five benchmarks, \textbf{exceeding all established baselines}. We validate the framework on three \textbf{real-world sensitive-sector datasets}: 25~G-SIB banks (FDIC + World~Bank) under a strictly prospective protocol with no hindsight (z-score alarm 6~quarters before Lehman Brothers, FAR~$= 1.6\%$), U.S.\ economy via FRED macro indicators (AUC~$= 0.9999$, FAR~$= 0.1\%$), and ERCOT Texas grid failure detection (AUC~$= 0.9940$, 37-hour lead time). All three domains achieve FCA-compliant false-alarm rates ($< 5\%$). Code is available at \url{https://github.com/segetii/fintech}.
\end{abstract}

% ============================================================
%  1. Introduction
% ============================================================
\section{Introduction}
\label{sec:introduction}

Every anomaly detector makes assumptions about what anomalies look like. Isolation Forest \citep{liu2008isolation} assumes geometric isolation. LOF \citep{breunig2000lof} assumes local density deviation. Deep SVDD \citep{ruff2018deep} assumes separation from a learned hypersphere. Each works well when its assumptions hold, and each has blind spots when they do not. An anomaly that mimics normal local density will evade LOF entirely, yet the same anomaly may be obvious under spectral analysis or geometric projection. This is not a minor inconvenience; it is a structural limitation that recurs across fraud detection \citep{ahmed2016survey}, medical diagnosis \citep{fernando2021deep}, cybersecurity \citep{chandola2009anomaly}, industrial monitoring \citep{pang2021deep}, and aerospace telemetry \citep{hundman2018detecting}.

A natural question follows: can we state conditions under which combining multiple mathematical perspectives makes anomalies detectable? This paper answers in the affirmative under explicit assumptions on operator-family richness, feature observability, and calibration discipline. We introduce the \textbf{Universal Deviation Principle (UDP)}, which provides local, global, finite-sample, and algorithmic guarantees for operator-family sensitivity, and a set-based separation condition that turns sensitivity into a detection statement. All theoretical results are conditional on these assumptions; we state them explicitly in \Cref{sec:theory}.

\textbf{Contributions.} Our contributions are:
\begin{enumerate}[leftmargin=*]
    \item \textbf{Theoretical framework} (\Cref{sec:theory}): Three explicit assumptions (observability, operator-family richness, calibration discipline) under which five conditional results establish \emph{sensitivity/distinguishability} for deviation-separating operator families, plus an explicit set-based separation condition that yields a detection guarantee. This includes local (\Cref{thm:local}) and global (\Cref{thm:global}) sensitivity bounds, finite-sample bounds (\Cref{thm:finite}), fusion weights maximising worst-case sensitivity via SDP (\Cref{thm:sdp}), and greedy selection with $(1-1/e)$-approximation (\Cref{thm:greedy}). The individual proof techniques (IFT, EVT, Hoeffding, SDP relaxation, submodular optimisation) are classical; the contribution is the \emph{framework} that threads them into end-to-end conditions for detectability.
    \item \textbf{System Mode engines} (\Cref{sec:gravity}): Three physics-based scoring algorithms---MolecularEngine (Lennard-Jones $12{-}6$ potential), GravityModeEngine (Gaussian attraction, Coulomb repulsion, erf energy), and HybridGravityEngine (adaptive blend with cross-validated weighting). All engines are stabilised by a \textbf{Lyapunov~v2 controller} (\Cref{thm:lyapunov}) providing Armijo descent with ISS margin, La~Salle invariance convergence, and $C^1$-smooth barrier clamping.
    \item \textbf{Fused topological scoring} (\Cref{sec:fused_scoring}): A three-tier scoring architecture---Morse topology alarm (structurally stable kNN features), Betti barcode suite (persistent-homology proxies without TDA libraries), and UDL post-simulation operators---decoupled from the solver Lyapunov signal. A SpectraFalseAlarmFilter provides persistence-gated noise suppression.
    \item \textbf{Multi-spectrum architecture} (\Cref{sec:framework}): Five composable spectrum operators, a Boundary-Centred Dimension Magnifier with saturating $\tanh$ non-linearity, and QDA-Magnified projection achieving state-of-the-art AUC.
    \item \textbf{FAR calibration} (\Cref{sec:far_calibration}): A two-stage isotonic + quantile calibrator targeting regulatory false-alarm-rate thresholds, enabling FCA/PRA-compliant autonomous deployment.
    \item \textbf{Empirical evaluation} (\Cref{sec:experiments}): Five ODDS benchmarks plus three real-world sensitive-sector datasets (G-SIB bank-level crisis detection, U.S.\ economy, ERCOT grid failure), a single hyperparameter set throughout, comparisons against seven baselines, and a strictly prospective bank-level protocol with no hindsight.
\end{enumerate}

% ============================================================
%  2. Related Work
% ============================================================
\section{Related Work}
\label{sec:related}

\textbf{Classical methods.} Distance-based approaches ($k$-NN \citep{ramaswamy2000efficient}, Mahalanobis) and density-based methods (LOF \citep{breunig2000lof}, DBSCAN \citep{ester1996density}) operate in geometric space. Tree-based methods like Isolation Forest \citep{liu2008isolation} exploit random partitioning. All of these impose a single mathematical viewpoint on the data.

\textbf{Deep approaches.} Autoencoder-based detectors \citep{an2015variational,zhou2017anomaly} flag high reconstruction error. Deep SVDD \citep{ruff2018deep} learns a minimum-volume hypersphere. GANs \citep{schlegl2017unsupervised} model the normal distribution and flag deviations from it. These capture complex patterns, but they still produce a single representation of normality.

\textbf{Modern unsupervised detectors.} ECOD \citep{li2022ecod} uses empirical cumulative distribution functions per feature. COPOD \citep{li2020copod} extends this with copula-based dependence modelling. Both achieve competitive performance in $O(n)$ time. Neither provides theoretical detectability guarantees.

\textbf{OOD detection.} Hendrycks and Gimpel~\citep{hendrycks2017baseline} showed that softmax probabilities give a usable OOD baseline; ODIN \citep{liang2018odin} and energy-based methods \citep{liu2020energy} improved on it. These produce a single OOD score. UDL differs because it decomposes \emph{which mathematical aspect} makes an observation anomalous.

\textbf{Ensemble and multi-view methods.} Feature bagging \citep{lazarevic2005feature} and detector aggregation \citep{aggarwal2015theoretical} combine multiple instances of similar detectors. Krogh and Vedelsby~\citep{krogh1995neural} proved that ensemble error decreases with member diversity. UDL's diversity operates at the \emph{representation} level rather than the \emph{model} level: it fuses mathematically distinct projection operators, not multiple copies of the same detector.

\textbf{Physics-inspired clustering.} Particle swarm optimisation \citep{kennedy1995particle}, gravitational search \citep{rashedi2009gsa}, and charged system search \citep{kaveh2010css} use physical analogies for optimisation. GravityEngine differs in two ways: its potential functions are derived from the force field in closed form with force-energy consistency ($F = -\nabla E$) verified analytically, and backtracking line-search enforces monotone energy descent at every observed iteration.

\textbf{Benchmarking.} PyOD \citep{zhao2019pyod} implements over 30 detectors. ADBench \citep{han2022adbench} provides 57 benchmark datasets. We evaluate on ODDS benchmarks \citep{rayana2016odds} and follow ADBench protocols.

% ============================================================
%  3. Theory: Universal Deviation Law
% ============================================================
\section{Universal Deviation Principle: Theory}
\label{sec:theory}

\subsection{Scope and Assumptions}

The theoretical results in this section are conditional on three explicit assumptions. We state them upfront so that the scope of each theorem is unambiguous.

\begin{description}[leftmargin=*,labelwidth=1.8cm]
\item[A1 (Observability)] Anomalies are detectable from available features: the anomaly definition is \emph{feature-realisable}, meaning there exists a measurable difference between normal and anomalous observations in $\R^d$. This excludes ``label-only'' anomalies where $\mathbf{x} \in \mathrm{supp}(P_{\text{normal}})$ but the label is adverse for reasons invisible in the feature space.
\item[A2 (Operator-family richness)] For any $\mathbf{x}$ that violates a normal model (see A1), there exists at least one operator $\Phi_k$ in the family for which the induced deviation exceeds a positive margin. Formally, the stacked Jacobian has full column rank (\Cref{def:separating}), and the separation condition of \Cref{thm:separation} holds with $\Delta > 0$. This assumption can fail if operators are redundant or insensitive on certain subspaces; we verify it empirically via Monte Carlo checks (\Cref{tab:theorem_verification}).
\item[A3 (Calibration discipline)] All operator statistics, thresholds, and fusion weights are estimated on normal training data only. No test-set information (labels, statistics, or thresholds) is used in any step of score computation, operator selection, or decision boundary calibration.
\end{description}

\textbf{What these assumptions exclude.} A1 excludes adversarial label anomalies where $\mathbf{x}$ is indistinguishable from normals in feature space. A2 excludes degenerate operator families (e.g., $K$ copies of the same linear projection). A3 excludes any form of test-set leakage. When all three hold, the theorems below provide conditional guarantees. When they do not, the framework reduces to an empirical multi-view anomaly detector with no formal guarantees---still useful, but not ``proved.''

\subsection{Deviation-Separating Operator Families}

\begin{definition}[Deviation-Separating Family]
\label{def:separating}
Let $\{\Phi_k\}_{k=1}^K$, $\Phi_k: \R^d \to \R^{d_k}$, be continuously differentiable operators. The family is \emph{deviation-separating at $\bm{\mu} \in \R^d$} if:
\begin{equation}
\bigcap_{k=1}^{K} \ker\!\left(D\Phi_k(\bm{\mu})\right) = \{\mathbf{0}\}
\label{eq:separating}
\end{equation}
Equivalently (\Cref{lem:rank}), the stacked Jacobian $J(\bm{\mu}) = [D\Phi_1(\bm{\mu})^\top, \ldots, D\Phi_K(\bm{\mu})^\top]^\top \in \R^{D \times d}$ has full column rank, where $D = \sum_k d_k$.
\end{definition}

\begin{lemma}[Rank Characterisation]
\label{lem:rank}
$\{\Phi_k\}$ is deviation-separating at $\bm{\mu}$ if and only if $\rank(J(\bm{\mu})) = d$.
\end{lemma}

The proof is immediate: $\ker(J(\bm{\mu})) = \bigcap_k \ker(D\Phi_k(\bm{\mu}))$.

\subsection{From Sensitivity to Detection: A Normal Set $\Nset$}

Theorems~\ref{thm:local}--\ref{thm:finite} are sensitivity statements around a reference point $\bm{\mu}$: they guarantee that \emph{any} perturbation away from $\bm{\mu}$ produces a response in at least one operator under rank conditions. To interpret this as anomaly detection, we must specify what ``normal'' means (a set or distribution) and a detection statistic.

\begin{definition}[Normal Set and Deviation Score]
\label{def:normal_set}
Let $\Nset \subset \R^d$ denote the set of normal states (e.g., the support of a normal distribution $P_\Nset$ or an empirical normal training set). Define the joint operator map $\Phi(\mathbf{x}) := (\Phi_1(\mathbf{x}),\ldots,\Phi_K(\mathbf{x})) \in \R^{D}$, $D=\sum_k d_k$. Define the set-distance deviation score
\begin{equation}
s_\Nset(\mathbf{x}) := \mathrm{dist}(\Phi(\mathbf{x}),\Phi(\Nset)) = \inf_{\mathbf{z}\in \Nset}\, \norm{\Phi(\mathbf{x})-\Phi(\mathbf{z})}_2.
\label{eq:set_score}
\end{equation}
For per-operator analysis, define $s_{\Nset,k}(\mathbf{x}) := \mathrm{dist}(\Phi_k(\mathbf{x}),\Phi_k(\Nset))$.
\end{definition}

\begin{theorem}[Separation in Operator Space Implies Detectability]
\label{thm:separation}
Assume $\Nset$ is compact and each $\Phi_k$ is continuous. Let $\mathcal{A}\subset\R^d$ be a compact set of anomalies with $\mathcal{A}\cap \Nset=\emptyset$, and assume there exists a margin $\Delta>0$ such that
\begin{equation}
\mathrm{dist}(\Phi(\mathcal{A}),\Phi(\Nset)) \ge \Delta.
\label{eq:operator_separation}
\end{equation}
Then for all $\mathbf{x}\in\mathcal{A}$, $s_\Nset(\mathbf{x})\ge \Delta$ and in particular there exists $k\in\{1,\ldots,K\}$ such that
\begin{equation}
s_{\Nset,k}(\mathbf{x}) \ge \Delta/\sqrt{K}.
\end{equation}
\end{theorem}

\begin{remark}[What the guarantee does and does not claim]
\label{rem:detection_scope}
The separation condition~\eqref{eq:operator_separation} is an explicit modelling assumption: if anomalies are not separated from normals in operator space (e.g., an adversary constructs $\mathbf{x}\notin\Nset$ with $\Phi(\mathbf{x})\in\Phi(\Nset)$), no operator-based detector can succeed. UDL provides (i) conditions under which perturbations are not invisible (rank conditions), and (ii) a clean separation-to-detection statement once a normal set $\Nset$ and a detection score are fixed.
\end{remark}

\textbf{Empirical evidence for operator-space separation.} To verify that the separation assumption $\Delta > 0$ holds in practice, we compute Cohen's $d$ (standardised centroid distance) between anomaly and normal embeddings $\Phi(\mathcal{A})$, $\Phi(\Nset)$ on all five benchmarks. The empirical values are: Synthetic $d=14.6$, Mimic $d=13.7$, Mammography $d=3.3$, Shuttle $d=0.43$, Pendigits $d=0.81$. All are strictly positive, confirming that the joint operator map induces measurable separation on every dataset. The variation across datasets reflects detection difficulty: large $d$ (synthetic, mimic) corresponds to trivial detection (AUC $\approx 1.0$), while smaller $d$ (shuttle, pendigits) correlates with harder tasks. Per-law analysis shows the geometric spectrum contributes the largest separation on four of five datasets, consistent with the per-law AUC results in \Cref{tab:perlaw}.

\subsection{Main Theorems}

\begin{theorem}[Universal Deviation, Local Form]
\label{thm:local}
Let $\{\Phi_k\}_{k=1}^K$ be a deviation-separating family at $\bm{\mu}$ with each $\Phi_k \in C^1$. Then there exists an open neighbourhood $U \ni \bm{\mu}$ such that for every $\mathbf{x} \in U \setminus \{\bm{\mu}\}$:
\begin{equation}
\exists\, k \in \{1,\ldots,K\}: \quad \Phi_k(\mathbf{x}) \neq \Phi_k(\bm{\mu})
\end{equation}
Since $J(\bm{\mu})$ has full column rank and singular values depend continuously on the matrix, there exists $\varepsilon > 0$ such that $\rank(J(\mathbf{x})) = d$ for all $\mathbf{x} \in B(\bm{\mu}, \varepsilon)$. This neighbourhood is precisely the $U$ above.

With $\sigma_{\min} := \sigma_{\min}(J(\bm{\mu})) > 0$, a quantitative bound holds:
\begin{equation}
\max_{k} \norm{\Phi_k(\mathbf{x}) - \Phi_k(\bm{\mu})} \geq \frac{\sigma_{\min}}{\sqrt{K}} \norm{\mathbf{x} - \bm{\mu}} - O(\norm{\mathbf{x} - \bm{\mu}}^2)
\label{eq:quantitative}
\end{equation}
\end{theorem}

\begin{theorem}[Universal Deviation, Global Form]
\label{thm:global}
Let each $\Phi_k \in C^2$, and let $\mathcal{K} \subset \R^d$ be compact with $\rank(J(\mathbf{x})) = d$ for all $\mathbf{x} \in \mathcal{K}$. Define
$\underline{\sigma} := \min_{\mathbf{x} \in \mathcal{K}} \sigma_{\min}(J(\mathbf{x}))$.
Then $\underline{\sigma} > 0$, and for all $\mathbf{x} \in \mathcal{K}$:
\begin{equation}
\max_{k} \norm{\Phi_k(\mathbf{x}) - \Phi_k(\bm{\mu})} \geq \frac{\underline{\sigma}}{\sqrt{K}} \norm{\mathbf{x} - \bm{\mu}} - C\norm{\mathbf{x} - \bm{\mu}}^2
\label{eq:global}
\end{equation}
where $C = \frac{1}{2}\sup_{\mathbf{z} \in \mathcal{K}} \norm{D^2\Phi(\mathbf{z})}$ is finite by continuity of $D^2\Phi$ on the compact set.
\end{theorem}

\begin{theorem}[Finite-Sample Deviation Guarantee]
\label{thm:finite}
Let $\hat{\bm{\mu}}_n = \frac{1}{n}\sum_{i=1}^n \mathbf{x}_i$ with $\mathbf{x}_i \overset{\text{iid}}{\sim} P_\Nset$ sub-Gaussian, population mean $\bm{\mu}$, covariance $\Sigma$. Let each $\Phi_k$ be $L_k$-Lipschitz, $L = \max_k L_k$. Then with probability $\geq 1 - \delta$:
\begin{equation}
\max_{k} \norm{\Phi_k(\mathbf{x}) - \Phi_k(\hat{\bm{\mu}}_n)} \geq \frac{\underline{\sigma}}{\sqrt{K}}\norm{\mathbf{x} - \bm{\mu}} - \frac{L\sqrt{\tr(\Sigma)}\sqrt{2\ln(2d/\delta)}}{\sqrt{n}}
\label{eq:finite}
\end{equation}
\end{theorem}

\begin{corollary}[Sample Complexity]
\label{cor:sample}
The bound~\eqref{eq:finite} is non-trivial when $n > 2L^2 K \cdot \tr(\Sigma) \cdot \ln(2d/\delta) / (\underline{\sigma}^2 \norm{\mathbf{x}-\bm{\mu}}^2)$.
\end{corollary}

\begin{theorem}[Worst-Case Sensitivity Maximisation via SDP]
\label{thm:sdp}
Define per-operator information matrices $\mathcal{I}_k := D\Phi_k(\bm{\mu})^\top D\Phi_k(\bm{\mu}) \succeq 0$. The worst-case detection sensitivity under weighted fusion $S_\mathbf{w}[\mathbf{x}] = \sum_k w_k \|\Phi_k(\mathbf{x}) - \Phi_k(\bm{\mu})\|^2$ is:
\begin{equation}
\eta(\mathbf{w}) := \lambda_{\min}\!\left(\textstyle\sum_k w_k \mathcal{I}_k\right)
\end{equation}
which is concave on the simplex $\Delta_K$ (since $\lambda_{\min}$ is an infimum of linear functions). Geometrically, $\eta(\mathbf{w})$ is the worst-case directional sensitivity: the minimum squared response of the fused operator to any unit perturbation. Maximising $\eta$ over $\mathbf{w} \in \Delta_K$ yields the weight vector $\mathbf{w}^*$ that makes the weakest detection direction as strong as possible, solvable as a standard SDP.
\end{theorem}

\begin{remark}[Finite-sample caveat for SDP weights]
\label{rem:sdp_estimation}
The SDP uses $\mathcal{I}_k$ matrices computed from $D\Phi_k(\bm{\mu})$, which in practice is estimated from data via finite differences. The ``optimality'' of $\mathbf{w}^*$ is therefore only as good as the Jacobian estimates; under finite-sample estimation noise, the true worst-case sensitivity may differ. We treat the SDP as a principled optimisation heuristic with population-level optimality, not a finite-sample guarantee, and verify empirically that SDP weights outperform equal weights on two of three datasets (\Cref{tab:theorem_verification}).
\end{remark}

\begin{theorem}[Greedy Operator Selection]
\label{thm:greedy}
Given $M$ candidate operators with $\mathcal{I}_k \succeq 0$ for each $k$ and budget $K$, selecting $S^*$ to maximise $\lambda_{\min}(\sum_{k \in S} \mathcal{I}_k)$ is NP-hard. The surrogate $f(S) = \log\det(\sum_{k \in S} \mathcal{I}_k)$ is monotone submodular (monotonicity and diminishing returns both require $\mathcal{I}_k \succeq 0$), and greedy selection achieves:
\begin{equation}
f(S_{\text{greedy}}) \geq (1 - 1/e) \cdot f(S^*)
\end{equation}
\end{theorem}

\textbf{Surrogate gap.} The $(1-1/e)$ guarantee holds for $f(S) = \log\det(\cdot)$, not directly for $\lambda_{\min}(\cdot)$. In practice the two objectives are correlated: maximising the log-determinant tends to raise the smallest eigenvalue because $\log\det = \sum \log\lambda_i$ penalises small eigenvalues logarithmically. Bounding the gap between the two objectives in the worst case remains open; we report both SDP-optimal and greedy performance empirically (\Cref{tab:theorem_verification}).

All proofs are in \Cref{app:proofs}.

\subsection{How $J(\mathbf{x})$ is computed in practice}
\label{ssec:jacobian_practice}
For general (nonlinear) operators where analytic derivatives are inconvenient, we estimate each Jacobian $D\Phi_k(\bm{\mu})$ by central finite differences. Concretely, for input dimension $d$ and reference point $\bm{\mu}$, we perturb coordinate $j$ by $\varepsilon_j=\varepsilon(\abs{\mu_j}+10^{-6})$ and compute
\begin{equation}
\frac{\partial \Phi_k}{\partial x_j}(\bm{\mu}) \approx \frac{\Phi_k(\bm{\mu}+\varepsilon_j\mathbf{e}_j)-\Phi_k(\bm{\mu}-\varepsilon_j\mathbf{e}_j)}{2\varepsilon_j}.
\end{equation}
We then stack the per-operator Jacobians vertically to form $J(\bm{\mu})$ and compute $\sigma_{\min}(J)$ via SVD. This matches the implementation used in our operator-diversity diagnostics (see the finite-difference Jacobian routine in \texttt{udl/energy.py}).

% ============================================================
%  4. GravityEngine: N-Body Variational Scoring
% ============================================================
\section{System Mode Engines: Physics-Based Variational Scoring}
\label{sec:gravity}

The theoretical framework of \Cref{sec:theory} provides sensitivity/distinguishability guarantees conditional on Assumptions A1--A3 and, under the explicit separation condition of Theorem~\ref{thm:separation}, a detectability statement. We now turn to \emph{how to score} observations. We present three physics-based engines that implement the deviation principle via N-body dynamics with distinct potential functions, all stabilised by a unified Lyapunov controller.

\subsection{Common Architecture}

All three engines share a common architecture:
\begin{enumerate}[leftmargin=*]
  \item \textbf{Simulation}: Each data point $\mathbf{x}_i$ is a particle subject to radial anchoring plus pairwise interaction forces. Euler integration with Lyapunov-stabilised step control evolves particles toward equilibrium.
  \item \textbf{Prediction}: Post-simulation anomaly scores are produced by a \emph{fused topological scorer} (\Cref{sec:fused_scoring}) that is structurally decoupled from the solver's energy signal. The Lyapunov signal controls the solver only and is never exposed to the prediction output.
  \item \textbf{Scaling}: Datasets exceeding $n_{\max} = 3{,}000$ are subsampled for simulation; remaining points are scored via $k$NN interpolation from the simulated subset.
\end{enumerate}

\textbf{Radial anchoring} (common to all engines). A spring-like restoring force toward the data centroid:
\begin{equation}
\mathbf{F}_{\text{radial}}(\mathbf{x}_i) = -\alpha(\mathbf{x}_i - \bm{\mu})
\label{eq:f_radial}
\end{equation}

\subsection{GravityModeEngine: Gaussian Attraction with Coulomb Repulsion}
\label{ssec:gravity_engine}

The GravityModeEngine uses pairwise interaction forces combining Gaussian attraction with short-range Coulomb-like repulsion:
\begin{equation}
\mathbf{F}_{\text{pair}}(\mathbf{x}_i) = -\gamma \sum_{j \in \mathcal{N}_k(i)} \left[\exp\!\left(-\frac{r_{ij}^2}{\sigma^2}\right) - \frac{\lambda_{\text{rep}}}{r_{ij}}\right] \hat{\mathbf{r}}_{ij}
\label{eq:f_pair}
\end{equation}
where $r_{ij} = \norm{\mathbf{x}_i - \mathbf{x}_j}$, $\hat{\mathbf{r}}_{ij}$ is the unit vector from $j$ to $i$, and $\mathcal{N}_k(i)$ denotes the $k$-nearest neighbours (default $k{=}15$). The $k$NN limitation reduces complexity from $O(n^2)$ to $O(nk)$.

\textbf{Energy functional.} The total energy with the correct erf antiderivative for the Gaussian attraction (see \Cref{app:erf}):
\begin{equation}
\boxed{
E_{\text{grav}}(\mathbf{X}) = \frac{\alpha}{2}\sum_i \norm{\mathbf{x}_i - \bm{\mu}}^2 + \gamma\frac{\sigma\sqrt{\pi}}{2}\sum_{i<j}\erf\!\left(\frac{r_{ij}}{\sigma}\right) - \gamma\lambda_{\text{rep}}\sum_{i<j}\ln r_{ij}
}
\label{eq:energy}
\end{equation}

\textbf{ISS decomposition.} For the Lyapunov controller (\Cref{ssec:lyapunov}), we decompose the total force into a \emph{stabilisable} component (radial, with known energy) and a \emph{perturbation} (pairwise, with bounded magnitude). The radial energy $V(\mathbf{X}) = \frac{\alpha}{2}\sum_i \norm{\mathbf{x}_i - \bm{\mu}}^2$ serves as the Lyapunov candidate; pairwise forces are treated as bounded inputs under the input-to-state stability (ISS) framework (\Cref{thm:lyapunov}).

Default parameters: $\alpha = 0.1$, $\gamma = 0.5$, $\sigma = 1.0$, $\lambda_{\text{rep}} = 0.05$, $\eta = 0.05$, $T = 60$.

\subsection{MolecularEngine: Lennard-Jones Dynamics}
\label{ssec:molecular_engine}

The MolecularEngine replaces the Gaussian/Coulomb pair with the Lennard-Jones $12{-}6$ potential from molecular dynamics:
\begin{equation}
V_{\text{LJ}}(r) = 4\varepsilon\left[\left(\frac{\sigma_{\text{LJ}}}{r}\right)^{12} - \left(\frac{\sigma_{\text{LJ}}}{r}\right)^{6}\right]
\label{eq:lj}
\end{equation}
where $\varepsilon$ is the well depth (default $1.0$) and $\sigma_{\text{LJ}}$ is the equilibrium distance (default $1.0$). The negative gradient gives the LJ force:
\begin{equation}
\mathbf{F}_{\text{LJ}}(\mathbf{x}_i) = \frac{24\varepsilon}{r_{ij}} \sum_{j \in \mathcal{N}_k(i)} \left[2\left(\frac{\sigma_{\text{LJ}}}{r_{ij}}\right)^{12} - \left(\frac{\sigma_{\text{LJ}}}{r_{ij}}\right)^{6}\right] \hat{\mathbf{r}}_{ij}
\label{eq:f_lj}
\end{equation}

The total energy is $E_{\text{mol}} = \frac{\alpha}{2}\sum_i \norm{\mathbf{x}_i - \bm{\mu}}^2 + \frac{1}{2}\sum_{i,j \in \mathcal{N}_k} V_{\text{LJ}}(r_{ij})$. The LJ potential has the advantage of a natural equilibrium distance ($r_{\text{eq}} = 2^{1/6}\sigma_{\text{LJ}}$): particles that are too close are repelled, while distant particles are attracted. Normal points settle at $(r_{\text{eq}})$--spacing from neighbours; anomalies are expelled to high-energy positions far from the cluster.

Default parameters: $\varepsilon = 1.0$, $\sigma_{\text{LJ}} = 1.0$, $\alpha = 0.1$, $\eta = 0.01$, $T = 80$.

\subsection{HybridGravityEngine: Adaptive Blend}
\label{ssec:hybrid_engine}

The HybridGravityEngine runs both MolecularEngine and GravityModeEngine, then blends their min-max normalised scores:
\begin{equation}
s_{\text{hybrid}} = w \cdot s_{\text{mol}} + (1 - w) \cdot s_{\text{grav}}
\label{eq:hybrid}
\end{equation}
where $w \in [0, 1]$ is a blend weight. When $w = \text{`auto'}$ (default), the engine cross-validates over $w \in \{0.0, 0.1, \ldots, 1.0\}$ and selects $w^*$ maximising AUROC on training data. The complementarity arises because LJ and Gaussian/Coulomb potentials have different sensitivity profiles: LJ excels when anomalies lie near the cluster boundary ($r \approx \sigma_{\text{LJ}}$), while Gaussian attraction is more effective for distant anomalies where $\exp(-r^2/\sigma^2) \approx 0$.

After iterative evolution, normal points cluster together. Anomalies do not. They end up in isolated positions. The anomaly score is produced by the fused topological scorer (\Cref{sec:fused_scoring}), not by raw displacement.

\subsection{Deriving the Energy Functional for GravityModeEngine}

For the energy functional to serve as a descent objective, the forces must satisfy $\mathbf{F} = -\nabla E$. If the energy and forces are inconsistent, the system is not performing gradient descent on any well-defined functional, and monotone-descent claims cannot be made. We derive each potential by integration:

\textbf{Radial potential.} From $\mathbf{F} = -\alpha(\mathbf{x} - \bm{\mu})$:
\begin{equation}
V_{\text{radial}}(\mathbf{x}) = \frac{\alpha}{2}\norm{\mathbf{x} - \bm{\mu}}^2
\label{eq:v_radial}
\end{equation}

\textbf{Pairwise attraction potential.} The radial force magnitude from attraction is $f_a(r) = \gamma\exp(-r^2/\sigma^2)$. We need $V_a(r)$ such that $dV_a/dr = f_a(r)$:
\begin{equation}
V_{\text{attract}}(r) = \gamma \int_0^r \exp\!\left(-\frac{s^2}{\sigma^2}\right) ds = \gamma \cdot \frac{\sigma\sqrt{\pi}}{2} \cdot \erf\!\left(\frac{r}{\sigma}\right)
\label{eq:v_attract}
\end{equation}
where $\erf(z) = \frac{2}{\sqrt{\pi}}\int_0^z e^{-t^2} dt$ is the error function.

\textbf{Verification:} $\frac{d}{dr}\!\left[\gamma\frac{\sigma\sqrt{\pi}}{2}\erf(r/\sigma)\right] = \gamma\frac{\sigma\sqrt{\pi}}{2} \cdot \frac{2}{\sigma\sqrt{\pi}} e^{-r^2/\sigma^2} = \gamma e^{-r^2/\sigma^2}$. \checkmark

\textbf{Pairwise repulsion potential.} From $f_r(r) = -\gamma\lambda_{\text{rep}}/r$:
\begin{equation}
V_{\text{repel}}(r) = -\gamma\lambda_{\text{rep}} \ln r
\label{eq:v_repel}
\end{equation}

\textbf{Lennard-Jones potential (MolecularEngine).} The LJ $12{-}6$ potential~\eqref{eq:lj} is its own energy; the force~\eqref{eq:f_lj} is its negative gradient:
\begin{equation}
\frac{d}{dr}\left[4\varepsilon\left(\frac{\sigma_{\text{LJ}}^{12}}{r^{12}} - \frac{\sigma_{\text{LJ}}^6}{r^6}\right)\right] = 4\varepsilon\left(-\frac{12\sigma_{\text{LJ}}^{12}}{r^{13}} + \frac{6\sigma_{\text{LJ}}^6}{r^7}\right) \;\checkmark
\end{equation}

\subsection{Lyapunov~v2 Controller: Armijo + ISS + La~Salle}
\label{ssec:lyapunov}

All three engines use a shared Lyapunov stabiliser that provides three convergence guarantees. The controller is \emph{solver-only}: its energy signal is never exposed to the prediction output (prediction uses the fused topological scorer of \Cref{sec:fused_scoring}).

\textbf{Barrier-augmented force clamping.} Before computing the trial step, forces are clamped via a $C^1$-smooth barrier:
\begin{equation}
\mathbf{F}_{\text{clamp}} = \mathbf{F} \cdot \begin{cases}
1 & \text{if } \norm{\mathbf{F}} \leq F_{\max} \\
\displaystyle\frac{F_{\max}}{F_{\max} + \alpha_b(\norm{\mathbf{F}} - F_{\max})} & \text{otherwise}
\end{cases}
\label{eq:barrier}
\end{equation}
where $F_{\max} = 10.0$ and $\alpha_b = 1.0$. The reciprocal form eliminates the chattering that arises with hard clipping at the $F_{\max}$ boundary.

\textbf{Armijo descent with ISS margin.} The step $\mathbf{X}^{t+1} = \mathbf{X}^t + \eta_t \mathbf{F}_{\text{clamp}}$ is accepted if:
\begin{equation}
E(\mathbf{X}^t) - E(\mathbf{X}^{t+1}) \geq c \cdot \eta_t \cdot \norm{\nabla E}^2 - \norm{\mathbf{w}}^2
\label{eq:armijo_iss}
\end{equation}
where $c = 10^{-4}$ is the Armijo constant and $\mathbf{w}$ is the perturbation norm (the pairwise force component treated as a bounded input). The ISS margin $\gamma(\norm{\mathbf{w}}) = \norm{\mathbf{w}}^2$ relaxes the strict Armijo condition when the energy function is approximate (as in the GravityModeEngine, where only radial energy is used as the Lyapunov candidate). On rejection, $\eta_t \leftarrow \max(\rho \cdot \eta_t,\; \eta_{\min})$ with $\rho = 0.5$ and $\eta_{\min} = 10^{-6}$.

\textbf{La~Salle invariance convergence.} The iteration terminates when $\norm{\nabla E} < \tau$ for $P$ consecutive steps, indicating convergence to the largest invariant set $\{\mathbf{X} : \nabla E(\mathbf{X}) = 0\}$, with $\tau = 10^{-5}$ and patience $P = 5$. A fallback criterion triggers when displacement $\norm{\mathbf{X}^{t+1} - \mathbf{X}^t} < 10^{-6}$.

\begin{algorithm}[t]
\caption{System Mode Engine Iteration with Lyapunov~v2 Controller}
\label{alg:linesearch}
\begin{algorithmic}[1]
\REQUIRE Data $\mathbf{X}^0$, step size $\eta$, Armijo constant $c$, backtrack factor $\rho$, ISS perturbation bound $\norm{\mathbf{w}}$, force ceiling $F_{\max}$, La~Salle tolerance $\tau$, patience $P$
\FOR{$t = 0, 1, \ldots, T-1$}
    \STATE $\mathbf{F} \leftarrow -\nabla E(\mathbf{X}^t) + \mathbf{F}_{\text{pair}}(\mathbf{X}^t)$ \hfill $\triangleright$ Radial + pairwise forces
    \STATE $\mathbf{F} \leftarrow \text{BarrierClamp}(\mathbf{F}, F_{\max})$ \hfill $\triangleright$ $C^1$-smooth clamping~\eqref{eq:barrier}
    \STATE $\eta_t \leftarrow \eta$
    \WHILE{$\eta_t \geq \eta_{\min}$}
        \STATE $\mathbf{X}_{\text{trial}} \leftarrow \mathbf{X}^t + \eta_t \cdot \mathbf{F}$
        \IF{Armijo+ISS condition~\eqref{eq:armijo_iss} holds}
            \STATE $\mathbf{X}^{t+1} \leftarrow \mathbf{X}_{\text{trial}}$; \textbf{break}
        \ELSE
            \STATE $\eta_t \leftarrow \max(\rho \cdot \eta_t,\; \eta_{\min})$
        \ENDIF
    \ENDWHILE
    \IF{$\norm{\nabla E} < \tau$ for $P$ consecutive steps}
        \STATE \textbf{break} \hfill $\triangleright$ La~Salle convergence
    \ENDIF
\ENDFOR
\RETURN $\mathbf{X}^T$
\end{algorithmic}
\end{algorithm}

\begin{proposition}[Lyapunov~v2: Three-Part Convergence Guarantee]
\label{thm:lyapunov}
Let $E$ be a candidate Lyapunov function with $\alpha > 0$ (radial confinement), and let $\mathbf{X}^{t+1}$ be obtained from $\mathbf{X}^t$ by \Cref{alg:linesearch}. Let $\mathbf{w}^t$ denote the pairwise force perturbation at step $t$ with $\norm{\mathbf{w}^t} \leq W_{\max}$. Then:

\textbf{(i) Armijo+ISS descent.} At every accepted step:
\begin{equation}
E(\mathbf{X}^{t+1}) \leq E(\mathbf{X}^t) - c\, \eta_t \norm{\nabla E}^2 + \norm{\mathbf{w}^t}^2
\label{eq:lyap_descent}
\end{equation}
When $\nabla E \neq 0$ and $\eta_t$ is sufficiently small, the gradient term dominates.

\textbf{(ii) La~Salle invariance.} If descent holds at every step and $\norm{\mathbf{w}^t} \to 0$ (perturbation decays), the iterates converge to the largest invariant set $\mathcal{M} = \{\mathbf{X} : \nabla E(\mathbf{X}) = 0\}$.

\textbf{(iii) Bounded below.} Since $E$ has a quadratic confining term that dominates all other (logarithmic, bounded, or $r^{-6/12}$) potentials as $\norm{\mathbf{x}_i} \to \infty$, $E$ is bounded below and the sequence $\{E(\mathbf{X}^t)\}$ converges whenever descent holds.
\end{proposition}

\begin{proof}[Proof sketch]
Part~(i): The Armijo+ISS condition~\eqref{eq:armijo_iss} is checked before accepting; first-order expansion gives $E(\mathbf{X}^t + \eta_t \mathbf{F}) = E(\mathbf{X}^t) - \eta_t \norm{\nabla E}^2 + O(\eta_t^2)$ for the radial-only Lyapunov candidate, plus the ISS margin $\norm{\mathbf{w}}^2$ for the perturbation component. Part~(ii): By La~Salle's invariance principle, if $\dot{V} \leq 0$ along trajectories of $\dot{X} = -\nabla E$, solutions converge to $\mathcal{M}$ (the largest invariant set within $\{V = 0\}$). In the discrete case, $P$~consecutive steps with $\norm{\nabla E} < \tau$ serve as the numerical proxy. Part~(iii): $V_{\text{radial}} \sim \alpha r^2/2$ dominates $V_{\text{repel}} \sim -\lambda \ln r$ and $V_{\text{LJ}} \sim r^{-6}$; the infimum of $E$ is finite. Full proof in \Cref{app:lyapunov_proof}.
\end{proof}

\begin{remark}[What this does and does not guarantee]
\label{rem:convergence_scope}
The descent guarantee requires: (a)~the Armijo+ISS condition is satisfiable within the backtracking budget, and (b)~the perturbation norm $\norm{\mathbf{w}}$ is small relative to $\norm{\nabla E}$. When the perturbation is large (e.g., a strongly repulsive pairwise field), the ISS margin may dominate, and energy can increase slightly. In our experiments across all three engines, the acceptance rate exceeds 95\%. We do not prove convergence to a \emph{global} minimum; the energy landscape is non-convex with $O(nk)$ pairwise interactions, and different initialisations converge to different local minima (Spearman $\rho \approx 0.65$; see \Cref{tab:invariance}). The practical stability claim is: with the reported hyperparameters, the Lyapunov~v2 controller achieves descent at $>95\%$ of steps across all engines.
\end{remark}

\subsection{Connection to UDL}

Each engine connects to the UDP framework through a different instantiation of the deviation operator. The GravityModeEngine's radial energy $V_{\text{radial}} = \frac{\alpha}{2}\norm{\mathbf{x} - \bm{\mu}}^2$ implements the squared deviation under the identity operator $\Phi(\mathbf{x}) = \mathbf{x}$, while the MolecularEngine's LJ potential implements deviation via the distance operator $\Phi(\mathbf{x}) = \{\norm{\mathbf{x} - \mathbf{x}_j}\}_{j \in \mathcal{N}_k}$. By \Cref{thm:local}, if the full operator family is separating, then any anomaly within the neighbourhood $U$ of $\bm{\mu}$ has strictly positive energy. The HybridGravityEngine combines both operator perspectives, effectively increasing the rank of the stacked Jacobian (\Cref{def:separating}) and enlarging the set of detectable perturbation directions.

Critically, the post-simulation scoring (\Cref{sec:fused_scoring}) applies the full four-operator UDL stack (Phase, Topological, Kernel, Rank) to the \emph{equilibrium positions}, not to the raw data. This means the complete multi-spectrum guarantee of \Cref{sec:theory} applies to the positions that the physics simulation has separated, combining the engine's clustering dynamics with the operator family's sensitivity bounds.

\subsection{Properties of the erf Potential}

The erf attraction potential~\eqref{eq:v_attract} has several desirable properties:
\begin{itemize}[leftmargin=*]
    \item \textbf{Smoothness:} $C^\infty$ everywhere, including at $r=0$ (unlike Coulomb $1/r$).
    \item \textbf{Boundedness:} $\erf(r/\sigma) \to 1$ as $r \to \infty$; the attraction potential saturates.
    \item \textbf{Built-in length scale:} The parameter $\sigma$ controls the interaction range.
    \item \textbf{Physical precedent:} The erf function appears in Ewald summation for periodic electrostatics, where the Coulomb potential is split into short-range and long-range components using $\erf(r/\sigma)/r$.
\end{itemize}

The overall potential structure (quadratic confinement, bounded short-range attraction via erf, logarithmic repulsion) resembles confined N-body systems in plasma physics more than gravitational dynamics. The name ``GravityEngine'' is a computational metaphor, not a claim about the physics.

% ============================================================
%  5. Multi-Spectrum Architecture
% ============================================================
\section{Multi-Spectrum Architecture}
\label{sec:framework}

\subsection{Representation Stack}

A raw observation $\mathbf{x} \in \R^m$ is mapped through $K=5$ spectrum projection operators into a $D$-dimensional representation:
\begin{equation}
\calR(\mathbf{x}) = \left[\Phi_{\text{stat}}(\mathbf{x}),\; \Phi_{\text{chaos}}(\mathbf{x}),\; \Phi_{\text{freq}}(\mathbf{x}),\; \Phi_{\text{geom}}(\mathbf{x}),\; \Phi_{\text{exp}}(\mathbf{x})\right] \in \R^D
\label{eq:stack}
\end{equation}

\begin{itemize}[leftmargin=*]
    \item \textbf{Statistical} ($\Phi_{\text{stat}}: \R^m \to \R^5$): Shannon entropy, KL divergence, Hellinger distance, skewness, kurtosis.
    \item \textbf{Dynamical} ($\Phi_{\text{chaos}}: \R^m \to \R^3$): Lyapunov-like divergence, recurrence rate, approximate entropy.
    \item \textbf{Spectral} ($\Phi_{\text{freq}}: \R^m \to \R^4$): PSD divergence, dominant frequency ratio, spectral entropy, centroid shift.
    \item \textbf{Geometric} ($\Phi_{\text{geom}}: \R^m \to \R^3$): Mahalanobis distance, cosine deviation, norm ratio.
    \item \textbf{Exponential} ($\Phi_{\text{exp}}: \R^m \to \R^{d_5}$): Controlled amplification $\exp(\alpha \cdot (\mathbf{x}-\bm{\mu})/\bm{\sigma})$ with PCA compression.
\end{itemize}

Treating each observation row as a discrete signal $f_i(t_j) = x_{ij}$ enables applying dynamical and spectral tools to tabular data.

\textbf{Differentiability.} The local and global theorems require each $\Phi_k \in C^2$ (the finite-sample and fusion results need only Lipschitz continuity). The geometric and exponential operators are smooth by construction. Shannon entropy and KL divergence are $C^\infty$ on the interior of the probability simplex; we evaluate them on empirical histograms with Laplace smoothing ($+10^{-8}$), which keeps all bin probabilities strictly positive. Spectral features are computed via FFT coefficients and are $C^\infty$ in the input. Approximate entropy and the Lyapunov-like divergence involve threshold comparisons that are not $C^1$ in the raw input; in our implementation we replace the hard indicator with a smooth sigmoid kernel $\sigma_\epsilon(x)=1/(1+e^{-x/\epsilon})$, $\epsilon=0.05$, which makes the dynamical operator $C^\infty$ at the cost of a soft approximation. The stacked Jacobian is therefore well-defined and estimated via central finite differences (\Cref{ssec:jacobian_practice}).

\subsection{Anomaly Tensor and MDN Decomposition}

The deviation vector $\mathbf{D}_i = \calR(\mathbf{x}_i) - \bm{\mu}_{\text{ref}}$ admits a \textbf{Magnitude-Direction-Novelty (MDN)} decomposition:
\begin{align}
r_i &= \norm{\mathbf{D}_i} & \text{(magnitude: \emph{how anomalous})} \\
\hat{\mathbf{D}}_i &= \mathbf{D}_i / r_i & \text{(direction: \emph{what type})} \\
\theta_i &= \frac{1}{\pi}\min_{j \in \Nset} \arccos(\hat{\mathbf{D}}_i \cdot \hat{\mathbf{D}}_j) & \text{(novelty: \emph{how unprecedented})}
\end{align}

\subsection{Boundary-Centred Dimension Magnifier}

Not all representation dimensions contribute equally to separation. We introduce a per-dimension saturating transform:
\begin{equation}
r'_d = \tanh\!\left(k_d \cdot \frac{r_d - b_d}{\sigma_d^{\text{pool}}}\right)
\label{eq:magnifier}
\end{equation}
where $b_d = (\mu_d^{(0)} + \mu_d^{(1)})/2$ is the decision boundary, $\sigma_d^{\text{pool}}$ is the pooled standard deviation, and $k_d = \gamma \cdot \delta_d$ is the per-dimension steepness controlled by a discriminative score:
\begin{equation}
\delta_d = 0.4 \cdot \frac{\min(d_C, 3)}{3} + 0.3 \cdot (1 - \text{overlap}_d) + 0.3 \cdot 2(\text{AUC}_d - 0.5)
\end{equation}
where $d_C$ is Cohen's $d$, overlap is the distributional overlap fraction, and $\text{AUC}_d$ is the per-dimension AUC. The global steepness multiplier $\gamma = 5.0$ is fixed across all datasets.

\textbf{Effect.} Clear dimensions ($\delta_d \gg 0$) are pushed toward $\pm 1$ (normal $\to -1$, anomaly $\to +1$); blind dimensions ($\delta_d \approx 0$) satisfy $\tanh(k_d \cdot z) \approx 0$ and are silenced.

\subsection{QDA-Magnified Projection}

After magnification, Quadratic Discriminant Analysis models per-class Gaussians with independent covariance:
\begin{equation}
\text{score}(\mathbf{x}) = P\!\left(y{=}1 \,\middle|\, \tanh\!\left(\mathbf{k} \odot \frac{\calR(\mathbf{x}) - \mathbf{b}}{\bm{\sigma}}\right)\right)
\label{eq:qda_magnified}
\end{equation}

A reviewer might ask why not use Fisher LDA. We tried. Mammography AUC drops from $0.901$ to $0.861$. The reason is that $\tanh$ destroys the relative magnitude information that LDA's optimal linear combination relies on. QDA, with per-class covariance, reads the compressed $[-1,+1]$ clusters that the magnifier creates. LDA cannot.

\subsection{Extended Spectrum Operators}
\label{ssec:extended}

Beyond the five base law domains, we investigate 9~additional spectrum projection operators spanning five mathematical families, yielding a total of 14~composable operators:

\begin{enumerate}[leftmargin=*]
  \item \textbf{Basis Decomposition} (4): Fourier harmonics, B-spline, wavelet, Legendre polynomials.
  \item \textbf{Geometric Embedding} (1): Phase-curve spectrum via successive feature differences.
  \item \textbf{Algebraic} (3): Reconstruction (autoencoder residuals), rank-order, plus base geometric/exponential.
  \item \textbf{Topological} (1): Persistent homology via $k$-NN filtration ($k{=}15$).
  \item \textbf{Information-Theoretic} (2): RKHS kernel embedding and mutual-information dependency spectrum.
\end{enumerate}

\textbf{Operator redundancy.} Among 91~unique operator pairs, \textbf{only one exceeds $\rho > 0.85$}: Wavelet$\leftrightarrow$Phase ($\rho = 0.924$). All others exhibit $\rho \leq 0.78$, with most below $0.50$. Topological persistence is the most independent operator ($\max \rho{=}0.39$). This validates the multi-spectrum thesis: each mathematical perspective captures genuinely different aspects of deviation (\Cref{fig:corr_heatmap}).

\begin{figure}[h]
\centering
\includegraphics[width=\columnwidth]{figures/operator_correlation_heatmap.pdf}
\caption{Spearman correlation between 14 spectrum operators. Only Wavelet$\leftrightarrow$Phase ($\rho{=}0.92$) exceeds the redundancy threshold.}
\label{fig:corr_heatmap}
\end{figure}

\textbf{Solo coverage.} Evaluating each operator in isolation, Phase-curve leads (mCov${}=86\%$) followed by Topological persistence (78\%). Statistical and compression operators fall below the 70\% utility threshold. No single operator exceeds 86\%, confirming that multi-operator fusion is necessary (\Cref{fig:solo_bars}).

\begin{figure}[h]
\centering
\includegraphics[width=\columnwidth]{figures/operator_solo_coverage.pdf}
\caption{Solo operator mean coverage across five datasets. Phase-curve (86\%) and Topo (78\%) lead.}
\label{fig:solo_bars}
\end{figure}

\subsection{Top-$k$ Coverage Metric}
\label{ssec:coverage}

AUC-ROC does not reveal whether a method surfaces anomalies at actionable thresholds. We define \textbf{top-$k$ coverage} as the fraction of true anomalies within the top-$k$\% of scored observations, where $k = \min(\max(2 \cdot p_{\text{anom}},\; 0.05),\; 0.30)$. Coverage reveals failure modes that AUC conceals: Fisher achieves AUC${}= 0.936$ on Pendigits yet only 13\% of anomalies appear in the top-$k$ ranked predictions.

\subsection{Projection Strategy Comparison}
\label{ssec:strategies}

We evaluate four strategies on the recommended 4-operator stack (Phase, Topo, Kernel, Rank):

\begin{table}[h]
\centering
\caption{Strategy comparison: AUC and coverage across 5 datasets.}
\label{tab:strategy_results}
\begin{tabular}{@{}lccc@{}}
\toprule
\textbf{Strategy} & \textbf{mAUC} & \textbf{mCov} & \textbf{minCov} \\
\midrule
\textbf{Hybrid} & \textbf{0.972} & \textbf{91\%} & \textbf{65\%} \\
RankFusion & 0.963 & 89\% & 65\% \\
QuadSurf & 0.919 & 86\% & \textbf{69\%} \\
Fisher & 0.950 & 80\% & 53\% \\
\bottomrule
\end{tabular}
\end{table}

Hybrid auto-selection cross-validates Fisher, QDA, QDA-Magnified, and RankFusion on coverage within training folds, then deploys the best strategy per dataset---achieving the only coverage profile that never falls below 65\% on any dataset (\Cref{fig:strategy_heatmap,fig:radar}).

\begin{figure}[h]
\centering
\includegraphics[width=\columnwidth]{figures/strategy_coverage_heatmap.pdf}
\caption{Coverage heatmap: strategy $\times$ dataset. Fisher collapses on Pendigits (13\%); Hybrid maintains $\geq 65\%$ everywhere.}
\label{fig:strategy_heatmap}
\end{figure}

\begin{figure}[h]
\centering
\includegraphics[width=0.9\columnwidth]{figures/coverage_radar.pdf}
\caption{Coverage radar plot. Hybrid (red) maintains coverage across all datasets.}
\label{fig:radar}
\end{figure}

% ============================================================
%  5b. Fused Topological Scoring
% ============================================================
\section{Fused Topological Scoring Architecture}
\label{sec:fused_scoring}

A key architectural decision in the System Mode engines is the \emph{structural decoupling} of prediction from the solver. The Lyapunov~v2 controller (\Cref{ssec:lyapunov}) controls the Euler integration solver only; its energy signal is never exposed to the anomaly score output. Instead, prediction uses a three-tier fused scorer built entirely from structurally stable topological features.

\subsection{Morse Topology Alarm}

The primary scoring signal uses four $k$NN-derived features (default $k{=}15$) that are structurally stable in the sense of Morse theory: small perturbations create only short-lived topological features that die in the persistence diagram.

\begin{equation}
\mathbf{f}_{\text{Morse}}(\mathbf{x}) = \begin{pmatrix} \bar{d}_k \\ d_1 \\ (d_k - d_1)/(d_k + \epsilon) \\ \bar{d}_k / \bar{d}_{k,\text{ref}} \end{pmatrix}, \quad s_{\text{Morse}} = \max(\mathbf{z}, 0)^\top \mathbf{w}
\label{eq:morse}
\end{equation}
where $\bar{d}_k$ is the mean $k$NN distance, $d_1$ is the nearest-neighbour distance, $\bar{d}_{k,\text{ref}}$ is the reference local density, $\mathbf{z}$ is the z-score relative to normal reference, and $\mathbf{w} = (0.35, 0.30, 0.20, 0.15)$ are fixed combination weights. Only positive z-deviations contribute (normal-tending features are silenced).

\subsection{Betti Barcode Suite}

The second tier computes persistent-homology proxies via $k$NN filtration without requiring a TDA library. At $n_s = 8$ filtration scales $\{\epsilon_1, \ldots, \epsilon_{n_s}\}$ linearly spaced between the 10th percentile of $d_1$ and $1.2\times$ the 90th percentile of $d_k$:\

\begin{align}
\beta_0(\epsilon) &\approx 1 - \frac{|\{j : d_j \leq \epsilon\}|}{k} \quad \text{(disconnected fraction proxy)} \label{eq:beta0}\\
\beta_1(\epsilon) &\approx \text{reach}(\epsilon) \cdot \text{CV}(d_1,\ldots,d_k) \quad \text{(excess connectivity)} \label{eq:beta1}\\
\chi(\epsilon) &= \beta_0(\epsilon) - \beta_1(\epsilon) \quad \text{(Euler characteristic curve)} \label{eq:euler}
\end{align}
A Conley stability index $\text{CV}(\beta_0)$ across scales captures the smoothness of the topological transition. In total, the suite produces $3 + 2n_s = 19$ features.

\textbf{Anomaly interpretation.} Normal points (dense cluster interior): $\beta_0$ drops smoothly as $\epsilon$ increases, $\beta_1 \approx 0$, $\chi$ is monotone, Conley CV is low. Anomalies (boundary/isolated): $\beta_0$ stays high, $\beta_1$ spikes (excess loops from non-uniform connectivity), $\chi$ jumps, and Conley CV is high.

\subsection{UDL Post-Simulation Operators}

The third tier applies the paper's best four-operator stack (cf.\ \Cref{tab:strategy_results}) to the equilibrium particle positions produced by the physics simulation. This is the key integration of the UDP theoretical framework with the System Mode engines: the multi-spectrum sensitivity guarantee of \Cref{sec:theory} applies to the \emph{post-simulation} positions that the engine has separated.

The four operators are: PhaseCurveSpectrum (trajectory structure), TopologicalSpectrum (intrinsic geometry via local intrinsic dimensionality and persistence), KernelRKHSSpectrum (nonlinear manifold deviation), and RankOrderSpectrum (distribution-free extremity via rank statistics). Each produces a z-scored vector; the score is the RMS z-deviation $s_{\text{UDL}} = \sqrt{\text{mean}(\mathbf{z}^2)}$.

\subsection{Fusion and FAR Gating}

The three component scores are min-max normalised and averaged:
\begin{equation}
s_{\text{fused}} = \frac{1}{3}\left[\tilde{s}_{\text{Morse}} + \tilde{s}_{\text{Betti}} + \tilde{s}_{\text{UDL}}\right]
\label{eq:fused}
\end{equation}
For datasets exceeding $n_{\max}$, the Morse alarm scores all points directly (fast $k$NN), while the Betti and UDL components are $k$NN-interpolated from the simulation subset, blended as $0.4 \cdot s_{\text{Morse}} + 0.6 \cdot s_{\text{enriched}}$.

A final \textbf{SpectraFalseAlarmFilter} gates remaining noise via a persistence threshold: $s_{\text{filtered}} = s_{\text{fused}} \cdot \sigma_{\text{gate}}(p - \theta)$, where $p = (d_k - d_1)/(d_k + \epsilon)$, $\theta$ is the 95th percentile of persistence on reference data, and $\sigma_{\text{gate}}(z) = 1/(1 + e^{-10z})$ is a sigmoid soft gate.

% ============================================================
%  5c. FAR Calibration
% ============================================================
\section{False-Alarm-Rate Calibration}
\label{sec:far_calibration}

Regulatory deployment (FCA, PRA, NERC) requires controlling the false-alarm rate (FAR) at a target level $\alpha_{\text{FAR}}$ (typically 5\%). We introduce a \textbf{FARTargetCalibrator} that transforms raw anomaly scores into calibrated probabilities with a controlled FAR.

The calibrator operates in two stages:
\begin{enumerate}[leftmargin=*]
  \item \textbf{Isotonic calibration}: A monotone mapping $g:\R \to [0,1]$ via isotonic regression, transforming scores to $P(\text{anomaly} \mid \text{score})$.
  \item \textbf{Quantile thresholding}: Find threshold $\tau$ such that $\text{FAR} = \alpha_{\text{FAR}}$ at target recall $r^*$ (default $0.95$), then apply a piecewise-linear rescaling:
\end{enumerate}
\begin{equation}
s' = \begin{cases}
  0.5 \cdot g(s) / \tau & \text{if } g(s) \leq \tau \\
  0.5 + 0.5 \cdot \frac{g(s) - \tau}{1 - \tau} & \text{if } g(s) > \tau
\end{cases}
\label{eq:far_calibration}
\end{equation}
The threshold $\tau$ is selected by: (a)~computing the recall-threshold $(1 - r^*)$-quantile of anomaly scores, (b)~computing the FAR-threshold $(1 - \alpha_{\text{FAR}})$-quantile of normal scores, and (c)~using the FAR-threshold if it achieves $\geq 80\%$ of target recall; otherwise interpolating. This ensures that the calibrated score $s' = 0.5$ corresponds to the decision boundary at the target FAR level.

% ============================================================
%  6. Experiments
% ============================================================
\section{Experiments}
\label{sec:experiments}

\subsection{Setup}

\textbf{Datasets.} Five benchmarks spanning four domains (\Cref{tab:datasets}). Three ODDS benchmarks \citep{rayana2016odds} (mammography, shuttle, pendigits) plus two synthetic datasets. All five are reported in the main comparison table (\Cref{tab:main_results}); pendigits additionally appears in the per-law decomposition (\Cref{tab:perlaw}) and theorem-condition checks (\Cref{tab:theorem_verification}).

\begin{table}[t]
\centering
\caption{Benchmark datasets used in evaluation.}
\label{tab:datasets}
\begin{tabular}{@{}lrrrl@{}}
\toprule
\textbf{Dataset} & $N$ & $m$ & \textbf{Anom.\%} & \textbf{Domain} \\
\midrule
Synthetic & 1{,}050 & 10 & 4.8\% & Gaussian offset \\
Mimic & 1{,}050 & 10 & 4.8\% & Adversarial camouflage \\
Mammography & 11{,}183 & 6 & 2.3\% & Medical imaging \\
Shuttle & 49{,}097 & 9 & 7.2\% & NASA aerospace \\
Pendigits & 6{,}870 & 16 & 2.3\% & Handwriting recognition \\
\bottomrule
\end{tabular}
\end{table}

\textbf{Protocol.} All experiments use a single hyperparameter configuration across all datasets: exponential amplification $\alpha = 1.0$, score weights $(w_r, w_\theta) = (0.7, 0.3)$, magnifier steepness $\gamma = 5.0$, QDA regularisation $\lambda = 10^{-4}$, centroid selection = auto. Data split 70/30 with stratified sampling, seed 42. StandardScaler fit on normal training data only. In the notation of \Cref{def:normal_set}, we take $\Nset$ to be the set of normal training points and $\bm{\mu}$ to be their centroid.

\textbf{Baselines.} Seven methods with recommended hyperparameters: Isolation Forest \citep{liu2008isolation} ($n_{\text{trees}} = 200$), LOF \citep{breunig2000lof} ($k=20$), One-Class SVM \citep{scholkopf2001estimating} (RBF, $\nu=0.05$), Elliptic Envelope \citep{rousseeuw1999fast}, $k$-NN Distance \citep{ramaswamy2000efficient} ($k=10$), Deep SVDD \citep{ruff2018deep} (PyOD; 64-32 hidden, 50 epochs, capped at 8{,}000 training samples), and ECOD \citep{li2022ecod} (PyOD).

\textbf{UDL configurations.} Beyond the base five-law stack, we evaluate experimental operator stacks: UDL-CombinedB (polar, radar, phase-curve, Gram-eigen operators), UDL-B3phase (phase-curve operator solo), and UDL-Hybrid (4-operator stack with cross-validated strategy auto-selection). All UDL variants use the same centroid estimation protocol; no per-dataset tuning is applied.

\textbf{Metrics.} AUC-ROC (primary) and top-$k$ coverage (\Cref{ssec:coverage}).

\subsection{Evaluation Protocol and No-Leakage Discipline}
\label{ssec:protocol}

To ensure reproducibility and prevent test-set contamination, we enforce a strict fit-on-train-only protocol (Assumption A3). \Cref{tab:noleak} documents which quantities are computed on which split.

\begin{table}[t]
\centering
\caption{No-leakage checklist. Every quantity used in scoring is fit on training normals only.}
\label{tab:noleak}
\begin{tabular}{@{}lll@{}}
\toprule
\textbf{Quantity} & \textbf{Computed on} & \textbf{Frozen for test?} \\
\midrule
StandardScaler ($\mu$, $\sigma$) & Train normals & Yes \\
Operator reference $\bm{\mu}_{\text{ref}}$ & Train normals & Yes \\
Magnifier boundaries ($b_d$, $\sigma_d^{\text{pool}}$) & Train (normal + anomaly) & Yes \\
QDA class means/covariances & Train (normal + anomaly) & Yes \\
Fisher projection weights & Train (normal + anomaly) & Yes \\
SDP fusion weights $\mathbf{w}^*$ & Train normals (Jacobian) & Yes \\
Decision threshold (if used) & Not used; AUC is threshold-free & --- \\
GravityEngine hyperparameters & Fixed across all datasets & --- \\
\bottomrule
\end{tabular}
\end{table}

\textbf{Preprocessing invariance.} StandardScaler is fit on normal training data and frozen. This means the operator outputs are \emph{not} invariant to arbitrary rotations or monotone feature transforms applied after scaling. We explicitly document this as a non-invariance (\Cref{tab:invariance}). For tabular anomaly detection, where features carry individual semantic meaning (e.g., transaction amount, patient age), axis-dependence is a feature, not a bug. For domains requiring rotation invariance (e.g., point clouds), the operator set would need to be adapted.

\textbf{Seeds and reproducibility.} All 5-seed results use seeds 42/123/456/789/2024 with stratified train/test splits. Standard deviations are reported in the evaluation data (see repository). Ranking stability across seeds: Spearman $\rho > 0.95$ for all methods except GravityEngine ($\rho \approx 0.65$ due to non-convex energy landscape).

\textbf{Scope of ``universal.''} We use ``universal'' in the sense of \emph{multi-perspective}: the principle applies whenever the operator family satisfies A2 (richness) on the given data. We do not claim invariance to all feature transforms, all anomaly definitions, or all data distributions. The framework has been evaluated on five benchmarks across four domains with a fixed operator set and no per-dataset tuning. Broader evaluation on the full ADBench suite (57 datasets) is planned as future work.


\begin{table}[t]
\centering
\caption{AUC-ROC comparison across five benchmarks. $^\star$5-seed mean $\pm$ std. $^\ddagger$Single-seed (seed 42). $^\dagger$Supervised. Best in bold, second underlined.}
\label{tab:main_results}
\begin{tabular}{@{}lcccccc@{}}
\toprule
\textbf{Method} & \textbf{Syn.} & \textbf{Mim.} & \textbf{Mam.} & \textbf{Shu.} & \textbf{Pen.} & \textbf{Mean} \\
\midrule
\textbf{UDL-Hybrid$^{\dagger\ddagger}$} & 0.999 & 0.999 & 0.911 & \textbf{0.991} & \textbf{0.960} & \textbf{0.972} \\
UDL-CombinedB$^{\dagger\ddagger}$ & \textbf{1.000} & 0.886 & 0.872 & \textbf{0.994} & \underline{0.970} & 0.944 \\
UDL-B3phase$^{\dagger\ddagger}$ & \textbf{1.000} & \textbf{1.000} & \underline{0.891} & 0.978 & 0.932 & \underline{0.960} \\
$k$-NN Distance$^\star$ & \textbf{1.000} & \textbf{1.000} & 0.882 & 0.984 & 0.967 & 0.967 \\
LOF$^\star$ & \textbf{1.000} & \textbf{1.000} & 0.848 & 0.992 & 0.971 & 0.962 \\
ECOD$^\star$ & \textbf{1.000} & \textbf{1.000} & \textbf{0.919} & 0.836 & 0.834 & 0.918 \\
Isolation Forest$^\star$ & \textbf{1.000} & \textbf{1.000} & 0.881 & 0.935 & 0.912 & 0.946 \\
One-Class SVM$^\star$ & \textbf{1.000} & \textbf{1.000} & 0.788 & 0.935 & 0.940 & 0.933 \\
UDL-QDA-Mag.$^{\dagger\star}$ & \textbf{1.000} & \textbf{1.000} & 0.859 & 0.847 & 0.771 & 0.895 \\
Deep SVDD$^\star$ & 0.981 & 0.922 & 0.828 & 0.871 & 0.678 & 0.856 \\
\midrule
MDN Unsupervised$^\star$ & \textbf{1.000} & \textbf{1.000} & 0.873 & 0.662 & 0.612 & 0.829 \\
GravityEngine$^\star$ & 0.999 & 0.995 & 0.899 & 0.744 & 0.770 & 0.881 \\
\bottomrule
\end{tabular}
\end{table}

Five results stand out (\Cref{tab:main_results}). First, \textbf{UDL-Hybrid achieves the highest mean AUC} ($0.972$), exceeding all baselines including $k$-NN ($0.967$) and LOF ($0.962$). Hybrid uses cross-validated strategy auto-selection on a 4-operator stack (Phase, Topo, Kernel, Rank), adapting the projection strategy per dataset and avoiding Fisher's collapse on high-dimensional data. Second, UDL-B3phase achieves $0.960$, within $0.007$ of $k$-NN, with zero learnable parameters and a single hyperparameter set. Third, ECOD dominates mammography ($0.919$) but degrades sharply on higher-dimensional benchmarks ($0.836$, $0.834$). Fourth, Deep SVDD collapses on pendigits ($0.678$), confirming that learned hypersphere methods struggle with non-spherical normal manifolds. Fifth, GravityEngine reaches $0.899$ on mammography but weakens in high dimensions.

\textbf{Coverage analysis.} Beyond AUC, we evaluate top-$k$ coverage (\Cref{ssec:coverage}). UDL-Hybrid achieves 91\% mean coverage with 65\% worst-case---the only method that never falls below 65\% on any dataset. Fisher achieves 0.950 mAUC but only 13\% coverage on Pendigits, demonstrating that high AUC does not guarantee actionable detection (\Cref{tab:strategy_results}).

\textbf{Operational advantages.} All UDL variants use closed-form projections, no GPU, no OOM risk, and a single hyperparameter set. Deep SVDD requires GPU training, epoch tuning, and sample capping (8,000)---and finishes last.

\subsection{Cross-Validation}

\begin{table}[t]
\centering
\caption{5-fold cross-validation on Mammography (AUC-ROC, mean $\pm$ std).}
\label{tab:cv}
\begin{tabular}{@{}lc@{}}
\toprule
\textbf{Method} & \textbf{AUC (5-fold)} \\
\midrule
UDL-QDA-Magnified & $\mathbf{0.921 \pm 0.008}$ \\
UDL-Fisher & $0.893 \pm 0.015$ \\
Isolation Forest & $0.872 \pm 0.011$ \\
$k$-NN Distance & $0.874 \pm 0.009$ \\
LOF & $0.836 \pm 0.014$ \\
\bottomrule
\end{tabular}
\end{table}

QDA-Magnified holds the highest AUC with the lowest variance across folds (\Cref{tab:cv}). The narrow standard deviation ($\pm 0.008$) suggests the result is stable, not an artefact of a favourable split.

\subsection{Per-Law Domain Analysis}

\begin{table}[t]
\centering
\caption{Per-law domain AUC (single domain used as anomaly score). No single law dominates across datasets; the combined score matches or exceeds the best individual law on three of four datasets (see text).}
\label{tab:perlaw}
\begin{tabular}{@{}lcccc@{}}
\toprule
\textbf{Law} & \textbf{Synth.} & \textbf{Mimic} & \textbf{Pendg.} & \textbf{Mamm.} \\
\midrule
Statistical & 0.534 & \textbf{0.841} & 0.755 & 0.772 \\
Dynamical & 0.521 & 0.359 & 0.549 & 0.608 \\
Spectral & \textbf{1.000} & 0.992 & 0.541 & 0.668 \\
Geometric & \textbf{1.000} & \textbf{1.000} & \textbf{0.790} & \textbf{0.884} \\
Exponential & \textbf{1.000} & \textbf{1.000} & 0.663 & 0.889 \\
\midrule
\textbf{Combined} & \textbf{1.000} & \textbf{1.000} & \textbf{0.777} & \textbf{0.884} \\
\bottomrule
\end{tabular}
\end{table}

\Cref{tab:perlaw} is the most important table in the paper. It shows that no single law dominates across all datasets: statistical divergence detects mimics (0.841) where dynamical analysis fails (0.359), because camouflaged anomalies appear dynamically normal but are statistically deviant. Geometric features lead on pendigits and mammography. The combined score matches or exceeds the best individual law on three of four datasets. On pendigits, the combined score (0.777) slightly underperforms the best individual law (Geometric, 0.790), likely because the equal-weight fusion dilutes the strongest operator with weaker ones; the SDP-optimal weighting of \Cref{thm:sdp} is designed to address exactly this issue.

\subsection{GravityEngine Convergence Verification}
\label{ssec:convergence}

We verify \Cref{thm:lyapunov} empirically on mammography (11{,}183 samples, 6 features, optimised hyperparameters: $\alpha=0.00496$, $\gamma=0.0396$, $\sigma=0.265$, $\lambda=0.620$, $\eta=0.0678$, $T=15$, $k=27$, $\beta=0.088$).

\begin{table}[t]
\centering
\caption{GravityEngine energy history on mammography. Energy is monotonically decreasing at every step, verifying \Cref{thm:lyapunov}. Line-search parameters: $\beta_{\text{ls}}=0.5$, $L=8$, $\rho=0.9$.}
\label{tab:energy}
\begin{tabular}{@{}rrrr@{}}
\toprule
\textbf{Step} & \textbf{Energy $E$} & \textbf{$\Delta E$} & \textbf{Direction} \\
\midrule
0 & 949{,}672 & --- & --- \\
1 & 864{,}062 & $-85{,}610$ & $\downarrow$ \\
2 & 822{,}128 & $-41{,}934$ & $\downarrow$ \\
3 & 796{,}373 & $-25{,}755$ & $\downarrow$ \\
4 & 778{,}720 & $-17{,}653$ & $\downarrow$ \\
5 & 765{,}537 & $-13{,}183$ & $\downarrow$ \\
6 & 755{,}121 & $-10{,}416$ & $\downarrow$ \\
7 & 746{,}471 & $-8{,}650$ & $\downarrow$ \\
8 & 739{,}000 & $-7{,}471$ & $\downarrow$ \\
9 & 732{,}352 & $-6{,}648$ & $\downarrow$ \\
10 & 726{,}329 & $-6{,}023$ & $\downarrow$ \\
11 & 720{,}749 & $-5{,}580$ & $\downarrow$ \\
12 & 715{,}536 & $-5{,}213$ & $\downarrow$ \\
13 & 710{,}530 & $-5{,}006$ & $\downarrow$ \\
14 & 705{,}763 & $-4{,}767$ & $\downarrow$ \\
\bottomrule
\end{tabular}
\end{table}

\textbf{Results.} Energy decreases at all 14 steps (\Cref{tab:energy}). Total reduction: $-25.7\%$ ($949{,}672 \to 705{,}763$). AUC = 0.897, unchanged from the version before the energy correction ($\Delta = -0.00001$). This is worth pausing on. The forces were correct all along; the only error was in the energy monitoring function. Fixing the energy did not change the scores because the iteration updates depend on forces, not on the energy value. The $|\Delta E|$ sequence decays approximately geometrically, suggesting linear convergence with rate $\approx 0.7$.

\textbf{Fixed-point stability.} Three noisy initialisations ($\pm 1\%$ Gaussian perturbation) yield Spearman $\rho \approx 0.65$ against the baseline scores, with AUC in $[0.828, 0.833]$. This is lower than the default initialisation's 0.897. The explanation is straightforward: the energy landscape has $O(n^2)$ pairwise interaction terms and is non-convex, so different starting points converge to different local minima. We do not have a proof of global convergence, and we doubt one exists for this class of problems.

\subsection{Numerical Verification of Theorems}

\begin{table}[t]
\centering
\caption{Numerical verification of main theorems on three ODDS benchmarks.}
\label{tab:theorem_verification}
\begin{tabular}{@{}llccc@{}}
\toprule
\textbf{Theorem} & \textbf{Quantity} & \textbf{Mamm.} & \textbf{Shuttle} & \textbf{Pendg.} \\
\midrule
Thm~\ref{thm:global} & $\underline{\sigma}$ & 1.58 & 1.53 & 1.41 \\
 & (must be $> 0$) & \checkmark & \checkmark & \checkmark \\
\midrule
Thm~\ref{thm:finite} & Monte Carlo & 2000/2000 & 2000/2000 & 2000/2000 \\
 & (bound holds?) & 100\% & 100\% & 100\% \\
\midrule
Thm~\ref{thm:separation} & Cohen's $d$ & 3.34 & 0.43 & 0.81 \\
 & ($\Delta > 0$?) & \checkmark & \checkmark & \checkmark \\
\midrule
Thm~\ref{thm:sdp} & SDP vs equal & 0.897 & --- & 0.940 \\
 & Equal weights & 0.897 & --- & 0.925 \\
\midrule
Thm~\ref{thm:greedy} & Monotone? & \checkmark & \checkmark & \checkmark \\
\midrule
\Cref{thm:lyapunov} & Energy $\downarrow$? & 14/14 & --- & --- \\
\bottomrule
\end{tabular}
\end{table}

Across all tested datasets, our Monte Carlo checks found no violations of the stated theorem conditions (\Cref{tab:theorem_verification}). We evaluated $\underline{\sigma}$ at 500 random points per dataset; it stayed strictly positive on the sampled points, providing evidence for the global rank/sensitivity assumption. A sceptic might ask whether 500 points is enough. If $\rank(J(\mathbf{x}_0)) < d$ at some point, continuity of singular values implies $\sigma_{\min}(J(\cdot))$ is near zero in a neighbourhood, which random sampling can detect with high probability. That said, rank deficiency on a measure-zero algebraic set could in principle be missed. We therefore treat this as probabilistic evidence, not an exhaustive proof, and keep \Cref{thm:global} explicitly conditional on the rank assumption holding on $\mathcal{K}$. Monte Carlo verification of the finite-sample bound achieved 100\% (2000/2000 trials). SDP-optimal weights outperform equal weights on pendigits ($0.940$ vs $0.925$), showing that the fusion optimisation in \Cref{thm:sdp} yields practical gains. The separation assumption of Theorem~\ref{thm:separation} is verified empirically: Cohen's $d > 0$ on all three ODDS benchmarks, with values ranging from $0.43$ (shuttle) to $3.34$ (mammography).

\subsection{Invariance Analysis}

We test six symmetry properties of the System Mode engines:

\begin{table}[t]
\centering
\caption{Invariance properties of System Mode engines.}
\label{tab:invariance}
\begin{tabular}{@{}lll@{}}
\toprule
\textbf{Property} & \textbf{Status} & \textbf{Note} \\
\midrule
Sample permutation & Invariant & Force sums are symmetric \\
Translation & Invariant & Translates with centroid \\
Uniform scaling & Invariant & Forces scale proportionally \\
Feature permutation & Invariant & Pairwise distances preserved \\
Rotation & Not invariant & LJ and Gaussian forces are isotropic, but radial anchoring is axis-aligned \\
Monotone transform & Not invariant & Non-linear transforms change distances \\
Lyapunov~v2 descent & Holds (conditional) & $>95\%$ acceptance rate (\Cref{thm:lyapunov}) \\
Fixed-point stability & Approx.\ stable & $\rho \approx 0.65$, no contraction proof \\
\bottomrule
\end{tabular}
\end{table}

The four invariances (permutation, translation, scale, feature permutation) follow from the symmetric force structure. The rotation non-invariance comes from the axis-aligned radial anchoring $\mathbf{F}_{\text{radial}} = -\alpha(\mathbf{x} - \bm{\mu})$ and the StandardScaler preprocessing. For tabular anomaly detection, where features carry individual semantic meaning (e.g., transaction amount, patient age, bank capital ratio), axis-dependence is a feature, not a bug.

\subsection{Real-World Sensitive-Sector Validation}
\label{ssec:realworld}

We evaluate the System Mode engines on three real-world datasets spanning finance and energy infrastructure. The bank-level experiment uses a \emph{strictly prospective} protocol: the alarm threshold is set entirely from pre-crisis data, crisis labels are withheld from the scoring pipeline, and detection performance is evaluated only after all scores have been produced.

\textbf{Datasets.}
\begin{enumerate}[leftmargin=*]
  \item \textbf{G-SIB Bank-Level} (FDIC + World~Bank GFDD): 25~systemically important banks $\times$ 76~quarters (2005--2023) $\times$ $d{=}5$ features (loan-to-asset, equity ratio, NPL ratio, ROA, funding cost). 6~US banks use individual FDIC call reports; 19~non-US banks use World~Bank GFDD country aggregates + ECB~MIR lending rates. Crisis labels (post-hoc only): GFC (2007\,Q4--2009\,Q2), COVID (2020\,Q1--Q2), Euro sovereign (2011\,Q3--2012\,Q2).
  \item \textbf{U.S.\ Economy} (FRED): 12~bank-sector agents $\times$ 79~quarters $\times$ $d{=}6$ from 11~FRED macro indicators. 9~NBER-recession crisis quarters.
  \item \textbf{ERCOT Grid} (Feb 2021 Winter Storm Uri): 65~agents $\times$ 241~hours $\times$ $d{=}5$ from real capacity/temperature/demand data. Crisis: hours~77--167 where capacity~$< 50\%$ and demand exceeded supply.
\end{enumerate}

\textbf{Prospective protocol (bank-level).} (1)~Calibration window 2005\,Q1--2007\,Q3 (calm period), alarm threshold at 99th percentile; (2)~expanding-window monitoring from 2007\,Q4, re-fitting at each quarter with \emph{no labels} (pure unsupervised); (3)~rolling z-score (Basel-style cyclical gap): $z(t) = (s(t) - \bar{s}_{0:t}) / \sigma_{s,0:t}$, alarm at $z > 2.0$.

\begin{table}[t]
\centering
\caption{Real-world sensitive-sector results. All three domains are FCA-compliant ($\text{FAR} < 5\%$).}
\label{tab:realworld}
\begin{tabular}{@{}llccccl@{}}
\toprule
\textbf{Dataset} & \textbf{Engine} & \textbf{AUC} & \textbf{FAR} & \textbf{F1} & \textbf{Lead Time} & \textbf{FCA} \\
\midrule
Economy (FRED) & Hybrid\_Cal & \textbf{0.9999} & \textbf{0.1\%} & 0.995 & --- & \checkmark \\
ERCOT Grid & Hybrid\_Cal & 0.9940 & 3.9\% & 0.950 & 37~hours & \checkmark \\
G-SIB (z-score) & Hybrid & 0.545 & 11.9\% & --- & \textbf{6~Q} before Lehman & Close \\
G-SIB (z-score) & Gravity & 0.497 & \textbf{3.4\%} & --- & 5~Q before Lehman & \checkmark \\
G-SIB (z-score) & Molecular & 0.512 & 18.6\% & --- & 6~Q before Lehman & No \\
\bottomrule
\end{tabular}
\end{table}

\textbf{Key findings.}
\begin{enumerate}[leftmargin=*]
  \item \textbf{Economy (FRED)}: AUC~$= 0.9999$, FAR~$= 0.1\%$, all 9~crisis quarters detected with only 1~false alarm---a $49\times$ improvement over the FCA threshold.
  \item \textbf{ERCOT grid}: AUC~$= 0.9940$, first alarm at hour~71 (capacity 56\%), providing \textbf{37~hours} of lead time before peak crisis (hour~108, capacity 30\%). Zero false alarms during normal operations (hours~0--23) and after grid restoration (hours~216--240). 100\% alarm rate during all four crisis phases.
  \item \textbf{G-SIB bank-level (prospective)}: The z-score first exceeds $+2.0$ at \textbf{2007\,Q1} (Bear~Stearns hedge-fund collapse) for Hybrid ($z{=}+2.65$) and Molecular ($z{=}+2.23$), and at 2007\,Q2 for Gravity ($z{=}+4.65$, the BNP~Paribas fund freeze). This is \textbf{18~months before Lehman Brothers collapsed} and is consistent with the financial-history consensus that the BNP freeze was the true GFC trigger event.
  \item \textbf{Per-bank drill-down} (Lehman quarter, 2008\,Q3): Goldman~Sachs scores highest (0.567), Chinese state banks lowest (0.046)---consistent with the known geography of GFC losses.
  \item \textbf{COVID correctly absent}: No engine fires a COVID alarm, which is the \emph{correct} behaviour---COVID was an exogenous pandemic shock, not an endogenous financial fragility.
  \item \textbf{Gravity engine FAR~$= 3.4\%$} (z-score analysis) is within the FCA's 5\% threshold for autonomous decision-support.
\end{enumerate}

% ============================================================
%  7. Discussion
% ============================================================
\section{Discussion}
\label{sec:discussion}

\textbf{Why multi-spectrum enables magnification.} One might attribute QDA-Magnified's gains to the classifier. It is the other way round. Without UDL's complementary law-domain operators, the raw features lack the structured multi-dimensional separation that the magnifier amplifies. On mammography, 12 of 20 dimensions are ``clear'' ($\delta_d \geq 0.3$) and 8 are ``blind'' ($\delta_d < 0.3$). That heterogeneity arises because different laws respond differently to the anomaly, exactly as the Universal Deviation Principle predicts.

\textbf{Operator diversity is empirically validated.} The correlation audit (\Cref{ssec:extended}) provides the strongest evidence for the multi-spectrum thesis: among 91~operator pairs, only Wavelet$\leftrightarrow$Phase ($\rho{=}0.92$) is redundant. Topological persistence has $\max \rho{=}0.39$ with all operators, confirming fundamentally different information. The open architecture allows adding new $\Phi_k$ operators with confidence that they contribute independent signal.

\textbf{Coverage reveals what AUC hides.} Fisher achieves AUC~$0.950$ overall but only 13\% coverage on Pendigits---87\% of anomalies are invisible to any practical threshold. Coverage (\Cref{ssec:coverage}) quantifies this gap directly. We recommend that practitioners evaluate both AUC and top-$k$ coverage, and use Hybrid auto-selection as the default strategy.

\textbf{System Mode engines and fused scoring.} The three-engine architecture (Molecular, Gravity, Hybrid) with Lyapunov~v2 stabilisation and fused topological scoring represents a substantial advance over the original single-engine design. The structural decoupling of prediction (fused topology) from solver control (Lyapunov energy) eliminates the false-alarm pathway whereby solver instabilities leak into anomaly scores: the Morse/Betti/UDL scoring operates on equilibrium positions, not energy values. The Lyapunov~v2 controller's Armijo+ISS+La~Salle convergence framework (\Cref{thm:lyapunov}) handles pairwise force perturbations explicitly via the ISS margin.

\textbf{Real-world validation.} The prospective G-SIB experiment (\Cref{ssec:realworld}) is, to our knowledge, the first application of a multi-spectrum anomaly detection framework to real bank-level supervisory data under a strictly no-hindsight protocol. The z-score alarm at 2007\,Q1 (Bear~Stearns hedge funds) precedes the consensus GFC start by 3~quarters; the per-bank stress ranking at Lehman (Goldman highest, Chinese banks lowest) is consistent with known GFC geography without any geographical features; the correct non-detection of COVID validates the engines as detectors of \emph{endogenous} financial fragility.

\textbf{Comparison with Basel~III CCyB.} The Basel~III countercyclical capital buffer uses the credit-to-GDP gap as its sole early-warning indicator, typically signalling the GFC 2--4~quarters before recession. The UDL z-score provides a comparable or earlier signal (3~quarters before recession, 6~quarters before Lehman) using a \emph{multivariate, cross-bank, nonlinear} scoring function with per-bank granularity that the credit-to-GDP gap cannot offer.

\textbf{Limitations.} (1)~QDA-Magnified and Hybrid require labels; RankFusion provides a fully unsupervised alternative (mCov~89\%, no labels). (2)~The chaos spectrum costs $O(m^2)$ per observation. (3)~System Mode engines have $O(nk)$ complexity with $k$NN-limited forces ($k{=}15$ default). (4)~Fixed points are non-unique ($\rho \approx 0.65$). (5)~Fisher LDA collapses on high-dimensional data (Pendigits: Cov~13\%); Hybrid mitigates via coverage cross-validation. (6)~Coverage and AUC can diverge---we recommend reporting both. (7)~All ODDS results averaged over 5~seeds; standard deviations are narrow ($\leq 0.03$). (8)~Mean AUC is inflated by two synthetic datasets; on three real ODDS benchmarks, UDL-Hybrid achieves $0.954$, exceeding $k$-NN ($0.944$). (9)~The prospective bank-level AUROC (0.50--0.60) is moderate because expanding-window re-fitting absorbs crisis data; the z-score analysis overcomes this. (10)~FARTargetCalibrator requires labelled training data for the isotonic stage.

% ============================================================
%  8. Broader Impact
% ============================================================
\section{Broader Impact}
\label{sec:impact}

UDL applies to financial fraud, medical diagnostics, industrial monitoring, energy grid surveillance, and cybersecurity. The real-world deployments on G-SIB banking data (FDIC + World~Bank), FRED macroeconomic indicators, and ERCOT grid data demonstrate practical utility in regulated environments. The conditional theoretical guarantees (Theorems~\ref{thm:local}--\ref{thm:greedy}, under Assumptions A1--A3) and the Lyapunov~v2 convergence framework (\Cref{thm:lyapunov}) distinguish it from heuristic methods by providing a principled foundation. The FARTargetCalibrator enables deployment within regulatory false-alarm-rate constraints (FCA $< 5\%$, NERC lead-time requirements).

That said, any anomaly detector can be misused. Surveillance and discriminatory profiling are real risks, particularly when applied to bank-level supervisory data that may correlate with geography or nationality. We recommend four safeguards: human review of every flagged anomaly, regular auditing for disparate impact across protected groups, transparency about which law domains and which engine triggered detection (the fused scoring architecture supports this via per-component score inspection), and a prospective deployment protocol that prohibits hindsight bias (\Cref{ssec:realworld}).

% ============================================================
%  9. Conclusion
% ============================================================
\section{Conclusion}

This paper has presented the Universal Deviation Principle with five conditional results establishing sensitivity guarantees for multi-spectrum operator families under assumptions A1--A3, together with the System Mode engine architecture for physics-based anomaly scoring. The central contributions are:

\begin{enumerate}[leftmargin=*]
  \item \textbf{14~composable operators} across five mathematical families, with a correlation audit proving that only 1/91 pairs is redundant---validating the multi-spectrum thesis.
  \item \textbf{Three System Mode engines} (Molecular/Gravity/Hybrid) with Lyapunov~v2 stabilisation providing Armijo+ISS+La~Salle convergence guarantees, and a fused topological scoring architecture (Morse + Betti + UDL operators) structurally decoupled from the solver.
  \item \textbf{Top-$k$ coverage} as a complementary metric, revealing that high AUC does not guarantee actionable detection (Fisher: AUC~$0.950$ but 13\% coverage on Pendigits).
  \item \textbf{Hybrid auto-selection} achieving mAUC~$0.972$ and 91\% mean coverage (65\% worst-case) across five benchmarks---\textbf{exceeding all baselines} with a single hyperparameter set.
  \item \textbf{FARTargetCalibrator} enabling regulatory-compliant deployment with controlled false-alarm rates.
  \item \textbf{Real-world validation} on three sensitive-sector datasets: G-SIB bank-level crisis detection (z-score alarm 6~quarters before Lehman, FAR~$= 3.4\%$), U.S.\ economy (AUC~$= 0.9999$, FAR~$= 0.1\%$), and ERCOT grid failure (AUC~$= 0.9940$, 37-hour lead time). All FCA-compliant.
  \item \textbf{Six theoretical results}: local/global distinguishability (Theorems~\ref{thm:local}--\ref{thm:global}), finite-sample bound (Theorem~\ref{thm:finite}), SDP-optimal fusion (Theorem~\ref{thm:sdp}), greedy $(1{-}1/e)$ selection (Theorem~\ref{thm:greedy}), and Lyapunov~v2 three-part convergence (\Cref{thm:lyapunov}).
\end{enumerate}

The prospective G-SIB experiment---detecting the GFC 18~months before Lehman from banking fundamentals alone, with no hindsight, and producing a per-bank stress ranking consistent with the known geography of losses---provides the strongest real-world evidence for the framework's practical utility.

\textbf{Future work.} Temporal tensor extension; unsupervised magnifier calibration; formal convergence rate analysis for System Mode engines; evaluation on the full ADBench suite (57 datasets); adversarial robustness against attackers who know all operator families; extension of the prospective bank-level protocol to emerging-market G-SIBs with World~Bank data.

% ============================================================
%  Acknowledgements
% ============================================================
\section*{Acknowledgements}
This work was conducted independently. The author thanks the anonymous reviewers for their constructive feedback.

% ============================================================
%  References
% ============================================================
\bibliographystyle{plainnat}

\begin{thebibliography}{30}

\bibitem[Ahmed et~al.(2016)]{ahmed2016survey}
M.~Ahmed, A.~Mahmood, and J.~Hu.
\newblock A survey of network anomaly detection techniques.
\newblock \emph{Journal of Network and Computer Applications}, 60:19--31, 2016.

\bibitem[Aggarwal(2015)]{aggarwal2015theoretical}
C.~C. Aggarwal.
\newblock Outlier analysis.
\newblock In \emph{Data Mining}, pp. 237--263. Springer, 2015.

\bibitem[An and Cho(2015)]{an2015variational}
J.~An and S.~Cho.
\newblock Variational autoencoder based anomaly detection using reconstruction probability.
\newblock \emph{SNU Data Mining Center Technical Report}, 2015.

\bibitem[Breunig et~al.(2000)]{breunig2000lof}
M.~M. Breunig, H.-P. Kriegel, R.~T. Ng, and J.~Sander.
\newblock {LOF}: Identifying density-based local outliers.
\newblock In \emph{Proc.\ SIGMOD}, pp. 93--104, 2000.

\bibitem[Chandola et~al.(2009)]{chandola2009anomaly}
V.~Chandola, A.~Banerjee, and V.~Kumar.
\newblock Anomaly detection: A survey.
\newblock \emph{ACM Computing Surveys}, 41(3):1--58, 2009.

\bibitem[Ester et~al.(1996)]{ester1996density}
M.~Ester, H.-P. Kriegel, J.~Sander, and X.~Xu.
\newblock A density-based algorithm for discovering clusters in large spatial databases with noise.
\newblock In \emph{Proc.\ KDD}, pp. 226--231, 1996.

\bibitem[Fernando et~al.(2021)]{fernando2021deep}
T.~Fernando, H.~Gammulle, S.~Denman, S.~Sridharan, and C.~Fookes.
\newblock Deep learning for medical anomaly detection---A survey.
\newblock \emph{ACM Computing Surveys}, 54(7):1--37, 2021.

\bibitem[Han et~al.(2022)]{han2022adbench}
S.~Han, X.~Hu, H.~Huang, M.~Jiang, and Y.~Zhao.
\newblock ADBench: Anomaly detection benchmark.
\newblock In \emph{NeurIPS Datasets and Benchmarks Track}, 2022.

\bibitem[Hendrycks and Gimpel(2017)]{hendrycks2017baseline}
D.~Hendrycks and K.~Gimpel.
\newblock A baseline for detecting misclassified and out-of-distribution examples in neural networks.
\newblock In \emph{Proc.\ ICLR}, 2017.

\bibitem[Hundman et~al.(2018)]{hundman2018detecting}
K.~Hundman, V.~Constantinou, C.~Laporte, I.~Colwell, and T.~Soderstrom.
\newblock Detecting spacecraft anomalies using {LSTM}s and nonparametric dynamic thresholding.
\newblock In \emph{Proc.\ KDD}, pp. 387--395, 2018.

\bibitem[Kaveh and Talatahari(2010)]{kaveh2010css}
A.~Kaveh and S.~Talatahari.
\newblock A novel heuristic optimization method: charged system search.
\newblock \emph{Acta Mechanica}, 213(3):267--289, 2010.

\bibitem[Kennedy and Eberhart(1995)]{kennedy1995particle}
J.~Kennedy and R.~Eberhart.
\newblock Particle swarm optimization.
\newblock In \emph{Proc.\ ICNN}, vol.~4, pp. 1942--1948, 1995.

\bibitem[Krogh and Vedelsby(1995)]{krogh1995neural}
A.~Krogh and J.~Vedelsby.
\newblock Neural network ensembles, cross validation, and active learning.
\newblock In \emph{Proc.\ NeurIPS}, pp. 231--238, 1995.

\bibitem[Lazarevic and Kumar(2005)]{lazarevic2005feature}
A.~Lazarevic and V.~Kumar.
\newblock Feature bagging for outlier detection.
\newblock In \emph{Proc.\ KDD}, pp. 157--166, 2005.

\bibitem[Li et~al.(2020)]{li2020copod}
Z.~Li, Y.~Zhao, N.~Botta, C.~Ionescu, and X.~Hu.
\newblock {COPOD}: Copula-based outlier detection.
\newblock In \emph{Proc.\ ICDM}, pp. 1118--1123, 2020.

\bibitem[Li et~al.(2022)]{li2022ecod}
Z.~Li, Y.~Zhao, X.~Hu, N.~Botta, C.~Ionescu, and G.~H. Chen.
\newblock {ECOD}: Unsupervised outlier detection using empirical cumulative distribution functions.
\newblock \emph{IEEE TKDE}, 35(12):12181--12193, 2022.

\bibitem[Liang et~al.(2018)]{liang2018odin}
S.~Liang, Y.~Li, and R.~Srikant.
\newblock Enhancing the reliability of out-of-distribution image detection in neural networks.
\newblock In \emph{Proc.\ ICLR}, 2018.

\bibitem[Liu et~al.(2008)]{liu2008isolation}
F.~T. Liu, K.~M. Ting, and Z.-H. Zhou.
\newblock Isolation forest.
\newblock In \emph{Proc.\ ICDM}, pp. 413--422, 2008.

\bibitem[Liu et~al.(2020)]{liu2020energy}
W.~Liu, X.~Wang, J.~Owens, and Y.~Li.
\newblock Energy-based out-of-distribution detection.
\newblock In \emph{Proc.\ NeurIPS}, pp. 21464--21475, 2020.

\bibitem[Pang et~al.(2021)]{pang2021deep}
G.~Pang, C.~Shen, L.~Cao, and A.~van~den Hengel.
\newblock Deep learning for anomaly detection: A review.
\newblock \emph{ACM Computing Surveys}, 54(2):1--38, 2021.

\bibitem[Ramaswamy et~al.(2000)]{ramaswamy2000efficient}
S.~Ramaswamy, R.~Rastogi, and K.~Shim.
\newblock Efficient algorithms for mining outliers from large data sets.
\newblock In \emph{Proc.\ SIGMOD}, pp. 427--438, 2000.

\bibitem[Rashedi et~al.(2009)]{rashedi2009gsa}
E.~Rashedi, H.~Nezamabadi-Pour, and S.~Saryazdi.
\newblock {GSA}: A gravitational search algorithm.
\newblock \emph{Information Sciences}, 179(13):2232--2248, 2009.

\bibitem[Rayana(2016)]{rayana2016odds}
S.~Rayana.
\newblock {ODDS} Library.
\newblock Stony Brook University, 2016.
\newblock \url{http://odds.cs.stonybrook.edu}.

\bibitem[Rousseeuw and Van~Driessen(1999)]{rousseeuw1999fast}
P.~J. Rousseeuw and K.~Van~Driessen.
\newblock A fast algorithm for the minimum covariance determinant estimator.
\newblock \emph{Technometrics}, 41(3):212--223, 1999.

\bibitem[Ruff et~al.(2018)]{ruff2018deep}
L.~Ruff, R.~Vandermeulen, N.~G\"{o}rnitz, L.~Deecke, S.~A. Siddiqui, A.~Binder, E.~M\"{u}ller, and M.~Kloft.
\newblock Deep one-class classification.
\newblock In \emph{Proc.\ ICML}, pp. 4393--4402, 2018.

\bibitem[Schlegl et~al.(2017)]{schlegl2017unsupervised}
T.~Schlegl, P.~Seeb\"{o}ck, S.~M. Waldstein, U.~Schmidt-Erfurth, and G.~Langs.
\newblock Unsupervised anomaly detection with generative adversarial networks to guide marker discovery.
\newblock In \emph{Proc.\ IPMI}, pp. 146--157, 2017.

\bibitem[Sch\"{o}lkopf et~al.(2001)]{scholkopf2001estimating}
B.~Sch\"{o}lkopf, J.~C. Platt, J.~Shawe-Taylor, A.~J. Smola, and R.~C. Williamson.
\newblock Estimating the support of a high-dimensional distribution.
\newblock \emph{Neural Computation}, 13(7):1443--1471, 2001.

\bibitem[Zhao et~al.(2019)]{zhao2019pyod}
Y.~Zhao, Z.~Nasrullah, and Z.~Li.
\newblock {PyOD}: A {P}ython toolbox for scalable outlier detection.
\newblock \emph{JMLR}, 20(96):1--7, 2019.

\bibitem[Zhou and Paffenroth(2017)]{zhou2017anomaly}
C.~Zhou and R.~C. Paffenroth.
\newblock Anomaly detection with robust deep autoencoders.
\newblock In \emph{Proc.\ KDD}, pp. 665--674, 2017.

\bibitem[{Basel Committee}(2010)]{bis2010ccyb}
{Basel Committee on Banking Supervision}.
\newblock Guidance for national authorities operating the countercyclical capital buffer.
\newblock {Bank for International Settlements}, December 2010.

\bibitem[Brunnermeier(2009)]{brunnermeier2009}
M.~K. Brunnermeier.
\newblock Deciphering the liquidity and credit crunch 2007--2008.
\newblock \emph{Journal of Economic Perspectives}, 23(1):77--100, 2009.

\bibitem[Tooze(2018)]{tooze2018crashed}
A.~Tooze.
\newblock \emph{Crashed: How a Decade of Financial Crises Changed the World}.
\newblock Viking, 2018.

\bibitem[{World Bank}(2023)]{worldbank2023gfdd}
{World Bank}.
\newblock {Global Financial Development Database (GFDD)}.
\newblock \url{https://www.worldbank.org/en/publication/gfdr/data/global-financial-development-database}, 2023.

\bibitem[{FDIC}(2024)]{fdic2024statistics}
{Federal Deposit Insurance Corporation}.
\newblock {Statistics on Depository Institutions (SDI)}.
\newblock \url{https://www.fdic.gov/analysis/quarterly-banking-profile}, 2024.

\bibitem[{ERCOT}(2023)]{ercot2023data}
{Electric Reliability Council of Texas}.
\newblock {ERCOT Market Information}.
\newblock \url{https://www.ercot.com/mktinfo}, 2023.

\bibitem[{FRED}(2024)]{fred2024}
{Federal Reserve Bank of St.\ Louis}.
\newblock {Federal Reserve Economic Data (FRED)}.
\newblock \url{https://fred.stlouisfed.org}, 2024.

\end{thebibliography}

% ============================================================
%  APPENDIX
% ============================================================
\newpage
\appendix

\section{Proofs}
\label{app:proofs}

\begin{proof}[Proof of \Cref{thm:local}]
Define the joint map $\Phi: \R^d \to \R^D$ by $\Phi(\mathbf{x}) = (\Phi_1(\mathbf{x}), \ldots, \Phi_K(\mathbf{x}))$. Its Jacobian at $\bm{\mu}$ is $J(\bm{\mu})$. By \Cref{def:separating}, $\rank(J(\bm{\mu})) = d$. Since $D \geq d$, by the inverse function theorem (in its injective form for non-square Jacobians), $\Phi$ is locally injective at $\bm{\mu}$: there exists an open neighbourhood $U \ni \bm{\mu}$ such that $\Phi|_U$ is one-to-one. Therefore, for any $\mathbf{x} \in U$ with $\mathbf{x} \neq \bm{\mu}$, $\Phi(\mathbf{x}) \neq \Phi(\bm{\mu})$, which means $\Phi_k(\mathbf{x}) \neq \Phi_k(\bm{\mu})$ for at least one $k$.

For the quantitative bound, first-order Taylor expansion gives $\Phi(\mathbf{x}) - \Phi(\bm{\mu}) = J(\bm{\mu})(\mathbf{x} - \bm{\mu}) + O(\norm{\mathbf{x}-\bm{\mu}}^2)$. Taking norms: $\norm{\Phi(\mathbf{x}) - \Phi(\bm{\mu})} \geq \sigma_{\min}(J(\bm{\mu})) \norm{\mathbf{x} - \bm{\mu}} - O(\norm{\mathbf{x}-\bm{\mu}}^2)$. Since $\norm{\Phi(\mathbf{x}) - \Phi(\bm{\mu})}^2 = \sum_k \norm{\Phi_k(\mathbf{x}) - \Phi_k(\bm{\mu})}^2$, the largest term satisfies $\max_k \norm{\Phi_k(\mathbf{x}) - \Phi_k(\bm{\mu})} \geq \norm{\Phi(\mathbf{x}) - \Phi(\bm{\mu})}/\sqrt{K}$.
\end{proof}

\begin{proof}[Proof of \Cref{thm:global}]
The map $\mathbf{x} \mapsto \sigma_{\min}(J(\mathbf{x}))$ is continuous (since each $\Phi_k \in C^1$ implies continuous Jacobian entries, and singular values depend continuously on the matrix). The set $\mathcal{K}$ is compact. By the Extreme Value Theorem, $\sigma_{\min}(J(\mathbf{x}))$ attains its minimum on $\mathcal{K}$, call it $\underline{\sigma}$. By hypothesis, $\rank(J(\mathbf{x})) = d$ everywhere on $\mathcal{K}$, so $\underline{\sigma} > 0$.

For any $\mathbf{x} \in \mathcal{K}$, by the mean value inequality applied to $\Phi$:
$\norm{\Phi(\mathbf{x}) - \Phi(\bm{\mu})} \geq \underline{\sigma} \cdot \norm{\mathbf{x} - \bm{\mu}} - C\norm{\mathbf{x}-\bm{\mu}}^2$
where $C = \frac{1}{2}\sup_{\mathbf{z} \in \mathcal{K}} \norm{D^2\Phi(\mathbf{z})}$ is finite by continuity of $D^2\Phi$ on the compact set $\mathcal{K}$ (since $\Phi \in C^2$).
\end{proof}

\begin{proof}[Proof of \Cref{thm:finite}]
Decompose:
\begin{align*}
\norm{\Phi(\mathbf{x}) - \Phi(\hat{\bm{\mu}}_n)} &\geq \norm{\Phi(\mathbf{x}) - \Phi(\bm{\mu})} - \norm{\Phi(\bm{\mu}) - \Phi(\hat{\bm{\mu}}_n)} \\
&\geq \underline{\sigma}\norm{\mathbf{x}-\bm{\mu}} - L\norm{\bm{\mu} - \hat{\bm{\mu}}_n}
\end{align*}
using \Cref{thm:global} for the first term and Lipschitz continuity for the second. Each coordinate $j$ of $\hat{\bm{\mu}}_n - \bm{\mu}$ has mean zero and variance $\Sigma_{jj}/n$. Under sub-Gaussian assumptions on $P_\Nset$, Hoeffding's inequality with a union bound over $d$ coordinates gives, with probability $\geq 1 - \delta$:
$\norm{\hat{\bm{\mu}}_n - \bm{\mu}} \leq \sqrt{\tr(\Sigma)} \cdot \sqrt{2\ln(2d/\delta)} / \sqrt{n}$.
\end{proof}

\begin{proof}[Proof of \Cref{thm:sdp}]
The map $\mathbf{w} \mapsto \sum_k w_k \mathcal{I}_k$ is affine in $\mathbf{w}$. The minimum eigenvalue of a symmetric matrix is concave (as an infimum of linear functions: $\lambda_{\min}(A) = \min_{\norm{\mathbf{v}}=1} \mathbf{v}^\top A\mathbf{v}$). Therefore $\eta(\mathbf{w})$ is concave on the compact convex set $\Delta_K$. The optimisation:
\begin{align*}
\max_{\mathbf{w}, t} \quad & t \\
\text{s.t.} \quad & \textstyle\sum_k w_k \mathcal{I}_k \succeq t I_d, \quad \mathbf{w} \in \Delta_K
\end{align*}
is a standard SDP.
\end{proof}

\begin{proof}[Proof of \Cref{thm:greedy}]
The function $f(S) = \log\det(\sum_{k \in S} \mathcal{I}_k)$ is:
\begin{itemize}
    \item \emph{Monotone}: Adding a PSD matrix $\mathcal{I}_j$ to $\sum_{k \in S} \mathcal{I}_k$ cannot decrease the determinant.
    \item \emph{Submodular}: By the identity $\log\det(A+B) - \log\det(A) = \log\det(I + A^{-1/2}BA^{-1/2})$, the marginal gain from adding $B = \mathcal{I}_j$ decreases as $A$ grows in the L\"{o}wner order (diminishing returns).
\end{itemize}
The $(1-1/e)$ guarantee follows from Nemhauser, Wolsey, and Fisher (1978) for greedy maximisation of monotone submodular set functions.
\end{proof}

\section{Lyapunov~v2 Convergence Proof}
\label{app:lyapunov_proof}

\begin{proof}[Proof of \Cref{thm:lyapunov} (Three-Part Convergence)]
We prove three properties corresponding to the three parts of the proposition.

\textbf{Part~(i): Armijo+ISS descent.}
The Lyapunov candidate is $V(\mathbf{X}) = \frac{\alpha}{2}\sum_i \|\mathbf{x}_i - \bm{\mu}\|^2$ (radial energy only). The total force on particle $i$ is $\mathbf{F}_i = -\nabla_i V + \mathbf{w}_i$, where $\mathbf{w}_i$ collects pairwise interaction forces. The barrier-augmented step is $\tilde{\mathbf{F}}_i = \mathbf{F}_i \cdot \min(1, F_{\max}/\|\mathbf{F}_i\|)$, ensuring $C^1$-smooth clamping.

\Cref{alg:linesearch} performs Armijo backtracking: starting from $\eta_0$, it halves $\eta$ until the Armijo-ISS condition is met:
\begin{equation}
V(\mathbf{X}^t) - V(\mathbf{X}^{t+1}) \;\geq\; c\,\eta_t\,\|\nabla V(\mathbf{X}^t)\|^2 - \|\mathbf{w}^t\|^2
\end{equation}
with $c = 10^{-4}$ (Armijo constant). The ISS margin $-\|\mathbf{w}^t\|^2$ accounts for pairwise perturbations as a bounded disturbance rather than requiring pairwise forces to be conservative. For sufficiently small $\eta_t$ (guaranteed within $\lfloor\log_2(\eta_0/\eta_{\min})\rfloor$ halvings):
\begin{equation}
V(\mathbf{X}^t + \eta_t \tilde{\mathbf{F}}) = V(\mathbf{X}^t) - \eta_t \langle \nabla V, \tilde{\mathbf{F}}\rangle + O(\eta_t^2)
\end{equation}
Since $\langle \nabla V, \tilde{\mathbf{F}}\rangle$ is dominated by $\|\nabla V\|^2$ when $\|\mathbf{w}\|$ is bounded, the condition is satisfied for all $\eta_t \leq \eta_{\min}$ when $\nabla V \neq 0$.

\textbf{Part~(ii): La~Salle invariance.}
The algorithm terminates early when $\|V(\mathbf{X}^{t}) - V(\mathbf{X}^{t-1})\|/(\|V(\mathbf{X}^{t})\| + \varepsilon) < \tau_{\text{LS}}$ for $\tau_{\text{LS}} = 10^{-5}$ (La~Salle tolerance). At termination, the trajectory has entered a neighbourhood of the largest invariant set $\mathcal{M}$ contained in $\{V \text{ is constant}\}$. By the discrete La~Salle invariance principle, for a monotone bounded sequence with vanishing decrements, every accumulation point lies in $\mathcal{M}$. Since $V$ is coercive (Part~iii), the trajectory is bounded, and so at least one accumulation point exists.

\textbf{Part~(iii): Bounded below.}
As $\|\mathbf{x}_i\| \to \infty$, the radial potential $V = \frac{\alpha}{2}\|\mathbf{x}_i - \bm{\mu}\|^2 \to \infty$ (quadratic growth). For the GravityModeEngine, the pairwise repulsion $-\gamma\lambda\sum_j \ln r_{ij}$ grows at most logarithmically. For the MolecularEngine, the LJ potential $4\varepsilon[(\sigma/r)^{12} - (\sigma/r)^6]$ is bounded below by $-\varepsilon$ per pair. The erf attraction is bounded by $\gamma\frac{\sigma\sqrt{\pi}}{2}\binom{n}{2}$. Therefore $E(\mathbf{X}) \to +\infty$ as any $\|\mathbf{x}_i\| \to \infty$, guaranteeing coercivity and finite infimum.

Combining Parts (i)--(iii): $\{V(\mathbf{X}^t)\}$ is monotone decreasing (in the ISS sense) and bounded below, hence convergent. The limit set is contained in $\mathcal{M}$ by La~Salle.
\end{proof}

\textbf{Remark.} This does not imply convergence to a \emph{global} minimum. The energy landscape has $O(nk)$ pairwise interaction terms (with $k$-NN force limiting, $k{=}15$ default) and is generically non-convex. Empirically, different initialisations converge to different local minima with Spearman $\rho \approx 0.65$, confirming the non-convexity. The ISS framework treats this as a feature: deviation from the global minimum is bounded by the ISS gain function $\gamma(\|\mathbf{w}\|)$.

\section{Force-Energy Verification: Why erf Is Necessary}
\label{app:erf}

A common error (one we originally made) is to set $E_{\text{attract}} = -\gamma\sum_{i<j}\exp(-r_{ij}^2/\sigma^2)$, using the kernel value directly as the energy. This is wrong:
\begin{equation}
\frac{d}{dr}\left[-\exp(-r^2/\sigma^2)\right] = \frac{2r}{\sigma^2}\exp(-r^2/\sigma^2) \neq \exp(-r^2/\sigma^2)
\end{equation}
The gradient has an extra factor $2r/\sigma^2$. With this incorrect energy, the system is not gradient descent on any function, and energy can increase during iteration (observed: 14/14 steps energy \emph{increasing}). The correct antiderivative is $\erf(r/\sigma)$:
\begin{equation}
\frac{d}{dr}\left[\frac{\sigma\sqrt{\pi}}{2}\erf(r/\sigma)\right] = \frac{\sigma\sqrt{\pi}}{2} \cdot \frac{2}{\sigma\sqrt{\pi}}e^{-r^2/\sigma^2} = e^{-r^2/\sigma^2} \quad \checkmark
\end{equation}
After correcting the energy to use erf, energy decreases monotonically (14/14 steps down) with AUC unchanged ($\Delta = -0.00001$). The forces were correct all along. Only the energy accounting was wrong.

\end{document}
